\documentclass[11pt,oneside]{article}

\usepackage[T1]{fontenc}
\usepackage{lmodern}
\usepackage{microtype}
\usepackage{geometry}
\geometry{margin=0.92in,headheight=15pt}
\usepackage{amsmath,amssymb,amsthm,mathtools}
\usepackage{mathrsfs}
\usepackage{booktabs,longtable,tabularx,array,multirow}
\newcolumntype{L}[1]{>{\raggedright\arraybackslash}p{#1}}
\usepackage{enumitem}
\usepackage{xcolor}
\usepackage{xurl}
\usepackage{hyperref}
\usepackage{listings}
\usepackage{fancyhdr}
\usepackage{titlesec}
\usepackage{needspace}

\definecolor{navy}{RGB}{25,66,112}
\definecolor{green}{RGB}{20,106,74}
\definecolor{violet}{RGB}{101,61,145}
\definecolor{brick}{RGB}{151,55,45}
\definecolor{gray}{RGB}{83,91,99}
\definecolor{lightgray}{RGB}{247,248,250}
\definecolor{rulegray}{RGB}{105,112,120}

\hypersetup{
  colorlinks=true,
  linkcolor=navy,
  citecolor=green,
  urlcolor=navy,
  pdftitle={Alaoglu--Erdos Conjecture: Dyadic Auxiliary Polynomials, Carlitz Multiplier Rigidity, and Rank-Sensitive Additive Targets},
  pdfauthor={Research continuation report},
  pdfsubject={New reductions and rigorous progress toward the Alaoglu--Erdos integer-exponent conjecture}
}
\urlstyle{same}

\pagestyle{fancy}
\fancyhf{}
\lhead{Alaoglu--Erd\H{o}s conjecture}
\rhead{Continuation report IV}
\cfoot{\thepage}
\renewcommand{\headrulewidth}{0.4pt}

\titleformat{\section}{\Large\bfseries}{\thesection}{0.75em}{}
\titleformat{\subsection}{\large\bfseries}{\thesubsection}{0.75em}{}
\titleformat{\subsubsection}{\normalsize\bfseries}{\thesubsubsection}{0.75em}{}

\setlist[itemize]{leftmargin=2em,itemsep=0.3em,topsep=0.4em}
\setlist[enumerate]{leftmargin=2.2em,itemsep=0.3em,topsep=0.4em}
\setlist[description]{leftmargin=3.25cm,labelsep=0.55cm,itemsep=0.35em,topsep=0.4em}
\setlength{\parindent}{0pt}
\setlength{\parskip}{0.58em}
\setlength{\emergencystretch}{3.5em}

\newtheorem{theorem}{Theorem}[section]
\newtheorem{proposition}[theorem]{Proposition}
\newtheorem{lemma}[theorem]{Lemma}
\newtheorem{corollary}[theorem]{Corollary}
\newtheorem{conjecture}[theorem]{Conjecture}
\newtheorem{question}[theorem]{Question}
\theoremstyle{definition}
\newtheorem{definition}[theorem]{Definition}
\newtheorem{target}[theorem]{Research target}
\newtheorem{experiment}[theorem]{Computational experiment}
\newtheorem{warning}[theorem]{Caution}
\theoremstyle{remark}
\newtheorem{remark}[theorem]{Remark}

\newcommand{\N}{\mathbb N}
\newcommand{\Nzero}{\mathbb N_0}
\newcommand{\Z}{\mathbb Z}
\newcommand{\Q}{\mathbb Q}
\newcommand{\R}{\mathbb R}
\newcommand{\C}{\mathbb C}
\newcommand{\Qbar}{\overline{\mathbb Q}}
\newcommand{\Fq}{\mathbb F_q}
\newcommand{\supp}{\operatorname{supp}}
\newcommand{\Span}{\operatorname{span}}
\newcommand{\rank}{\operatorname{rank}}
\newcommand{\Res}{\operatorname{Res}}
\newcommand{\ord}{\operatorname{ord}}
\newcommand{\Log}{\operatorname{Log}}
\newcommand{\ee}{\mathrm e}
\newcommand{\cl}{\operatorname{cl}}
\newcommand{\AEconj}{Alaoglu--Erd\H{o}s}
\newcommand{\tht}{\theta}
\newcommand{\Nm}{\mathcal N}
\newcommand{\cT}{\mathcal T}
\newcommand{\cP}{\mathcal P}
\newcommand{\cC}{\mathcal C}
\newcommand{\cE}{\mathcal E}
\newcommand{\cG}{\mathcal G}
\newcommand{\cI}{\mathcal I}
\newcommand{\cL}{\mathcal L}
\newcommand{\cO}{\mathcal O}
\newcommand{\cR}{\mathcal R}
\newcommand{\cS}{\mathcal S}
\newcommand{\abs}[1]{\left|#1\right|}
\newcommand{\norm}[1]{\left\lVert#1\right\rVert}
\newcommand{\floor}[1]{\left\lfloor#1\right\rfloor}
\newcommand{\ceil}[1]{\left\lceil#1\right\rceil}
\newcommand{\statusproved}{\textcolor{green}{\textbf{[proved here]}}}
\newcommand{\statusimported}{\textcolor{gray}{\textbf{[imported theorem]}}}
\newcommand{\statuscomputed}{\textcolor{violet}{\textbf{[computed here]}}}
\newcommand{\statustarget}{\textcolor{navy}{\textbf{[open target]}}}
\newcommand{\statusnegative}{\textcolor{brick}{\textbf{[negative audit]}}}

\newenvironment{statusbox}[1]
{\begin{center}\begin{minipage}{0.95\textwidth}\raggedright\hrule
\vspace{0.45em}\textbf{Status:} #1\par\vspace{0.35em}}
{\vspace{0.35em}\hrule\end{minipage}\end{center}}

\newenvironment{resultbox}
{\begin{center}\begin{minipage}{0.95\textwidth}\hrule\vspace{0.5em}}
{\vspace{0.45em}\hrule\end{minipage}\end{center}}

\lstdefinestyle{pythonstyle}{
  language=Python,
  basicstyle=\ttfamily\footnotesize,
  backgroundcolor=\color{lightgray},
  frame=single,
  rulecolor=\color{rulegray},
  numbers=left,
  numberstyle=\tiny\color{rulegray},
  numbersep=8pt,
  showstringspaces=false,
  breaklines=true,
  columns=fullflexible,
  keepspaces=true,
  tabsize=4,
  captionpos=b
}
\lstdefinestyle{pseudostyle}{
  basicstyle=\ttfamily\footnotesize,
  backgroundcolor=\color{lightgray},
  frame=single,
  rulecolor=\color{rulegray},
  showstringspaces=false,
  breaklines=true,
  columns=fullflexible,
  keepspaces=true,
  tabsize=2
}

\begin{document}
\pagenumbering{gobble}

\begin{titlepage}
\centering
\vspace*{0.6cm}
{\Huge\bfseries The Alaoglu--Erd\H{o}s Conjecture\par}
\vspace{0.45cm}
{\LARGE\bfseries Dyadic Auxiliary Polynomials,\par}
\vspace{0.12cm}
{\LARGE\bfseries Carlitz Multiplier Rigidity, and\par}
\vspace{0.12cm}
{\LARGE\bfseries Rank-Sensitive Additive Targets\par}
\vspace{0.75cm}
{\Large Continuation Report IV\par}
\vspace{0.25cm}
{\large 23 August 2026\par}
\vfill
\begin{minipage}{0.92\textwidth}
\small
\textbf{Bottom line.}
No proof of the conjecture is claimed.  The preceding report was audited structurally---it
contains 23,520 lines and about 123,000 words---and the present document deliberately avoids
re-running its determinant, Kummer, carry-language, interpolation, compact-orbit, special-function,
and finite-kernel developments.  Instead it develops three directions absent from that report.

First, the identity $2^{\log_2 3}=3$ puts every dyadic block of the curve
$Y=X^{\log_2 3}$ \emph{exactly} into the fixed-denominator framework of the June 2026
Brisebarre--Hanrot algorithm.  Combining this normalization with an elementary Chebyshev-system
zero theorem gives a new support-pressure criterion: if the conjecture is false, then every
integer polynomial that is uniformly subunit on a whole normalized dyadic block must have
linearly many nonzero monomials.  More generally, the total monomial budget of any certified
cover must be at least linear in the block index.  Thus a cover with total support $o(L)$ on the
$L$th block would prove the conjecture.  This criterion distinguishes a hypothetical irrational
counterexample from all integral controls.

Second, Papanikolas's algebraic-independence theorem for Carlitz logarithms implies a rigorous
positive-characteristic analogue: if two $k$-linearly independent Carlitz logarithms remain
Carlitz logarithms after multiplication by the same scalar $c$, then $c$ lies in the rational
function field $k$.  Equivalently, the quadratic determinant relation factors through linear
relations.  Estienne's 2026 Mahler-method proof identifies the missing characteristic-zero
ingredient with unusual precision: a nontrivial difference-operator module that promotes
linear independence of logarithms to algebraic independence.

Third, Harrison's July 2026 theorem on additive relations in irrational powers is recast for the
multiplicative two-generator grid.  Its fixed-arity estimates remain compatible with the inherited
carry fibers; the needed theorem must be rank-sensitive in logarithmic coordinates and uniform
through arity proportional to logarithmic height.  The report states that target quantitatively
and records why simpler congruence-continuity and stationary-scaling ideas fail even on integral
controls.
\end{minipage}
\vfill
{\small All statements are tagged as proved here, imported, computed, open targets, or negative audits.}
\end{titlepage}

\clearpage
\pagenumbering{roman}
\tableofcontents
\clearpage
\pagenumbering{arabic}

\section{Scope, audit discipline, and status ledger}

\subsection{What is new, and what is intentionally not repeated}

The input report already develops the exact conditional semigroup of solutions, the canonical
primitive odd core, the four-exponentials barrier, rational-function closure, valuation lattices,
Smith profiles, arbitrary-node determinants, generalized Vandermonde bounds, $p$-adic logarithms,
Kummer support densities, carry languages, o-minimal counting, Sturmian dynamics, Mahler and
special-function transfer, entire interpolation, cyclotomic and factorial cocycles, a finite
verified kernel, and a large Lean-oriented dependency map.  Rephrasing those sections would add
volume but not information.

The present report therefore imposes the following non-duplication rule.

\begin{quote}
A topic from Report III is used only as an explicitly named input.  A section is included here
only if it supplies a new theorem, a new exact reduction, a new quantitative threshold, a new
algorithmic specialization, or a newly falsified route.
\end{quote}

The principal new observation is not merely that the curve $Y=X^{\tht}$ is analytic.  It is that
its two natural scales are arithmetically synchronized:
\[
  (2^L)^{\tht}=3^L,\qquad \tht:=\frac{\log 3}{\log 2}.
\]
This removes the scale-dependent function that normally appears when a power curve is normalized
to $[1,2]$.  The function remains exactly $z\mapsto z^{\tht}$, while the two fixed denominators are
$2^L$ and $3^L$.  That feature is what makes the 2026 near-curve algorithm unusually well matched
to this problem.

\subsection{Status convention}

\begin{longtable}{L{0.18\textwidth}L{0.75\textwidth}}
\toprule
Tag & Meaning \\
\midrule
\endfirsthead
\toprule
Tag & Meaning \\
\midrule
\endhead
\statusproved & A complete proof is supplied in this report, using only stated inputs. \\
\statusimported & The statement is taken from a cited paper and is not reproved here. \\
\statuscomputed & Numerical values were recomputed by the script in Appendix~\ref{app:script}. \\
\statustarget & A precise sufficient condition or research problem; it is not asserted to hold. \\
\statusnegative & A tempting route is disproved, usually by an integral control example. \\
\bottomrule
\end{longtable}

\subsection{Ledger of the actual advances}

\begin{longtable}{L{0.28\textwidth}L{0.18\textwidth}L{0.44\textwidth}}
\toprule
Item & Status & Content \\
\midrule
\endfirsthead
\toprule
Item & Status & Content \\
\midrule
\endhead
Exact dyadic normalization & \statusproved & The $L$th dyadic block is exactly Problem 1.1 of Brisebarre--Hanrot with fixed $f(z)=z^{\tht}$, $u=2^L$, and $v=3^L$. \\
Counterexample block pressure & \statusproved & Any single noninteger solution $x$ creates at least $\ceil{(L+1)/x}$ distinct integer points in every sufficiently indexed dyadic block. \\
Fewnomial zero budget & \statusproved & A polynomial with $T$ nonzero monomials meets the positive curve $Y=X^{\tht}$ in at most $T-1$ points. \\
Sparse-certificate criterion & \statusproved & A uniformly subunit cover whose total monomial budget is $o(L)$ proves the conjecture. \\
Uniform Bernstein ellipse & \statusproved & $z^{\tht}$ is analytic on every normalized Bernstein ellipse with $1<\rho<3+2\sqrt2$, with explicit uniform norm bounds. \\
Near-curve parameter translation & \statuscomputed & The 2026 complexity exponents are translated into dyadic-block variables and tabulated. \\
Carlitz multiplier rigidity & \statusproved & Papanikolas's theorem implies that a common multiplier preserving two independent algebraic Carlitz logarithms lies in $\Fq(\vartheta)$. \\
Quadratic relation-ideal target & \statustarget & The degree-two characteristic-zero analogue would settle the original conjecture immediately. \\
Toric semi-invariant classification & \statusproved & A polynomial satisfying $P(2X,3Y)=cP(X,Y)$ is a monomial.  Pure scaling therefore cannot be the missing operator. \\
Log-GAP growing-arity target & \statustarget & Harrison's theorem is reformulated at the rank and arity actually forced by a counterexample. \\
Congruence-continuity idea & \statusnegative & A same-prime kernel inclusion already fails for the valid control exponent $x=2$ modulo $7$. \\
\bottomrule
\end{longtable}

\subsection{External 2026 context}

Claude Fable 5 was announced on 9 June 2026~\cite{AnthropicFable2026}.  Publicly checkable
sources support several parts of the competitive framing, but at different evidentiary levels.  An
explicit three-variable counterexample attributed to Fable refutes the general Jacobian conjecture
in dimensions at least three, while the plane case remains open~\cite{JacobianReport2026}.  Epoch
AI has provisionally marked the order-$668$ Hadamard challenge as solved by AI and records that the
reported construction was credited to three humans and Claude~\cite{EpochHadamard2026}; because
$668$ was the smallest of the twelve historically unresolved admissible orders below $2000$, the
announced family would clear that old list.  Independently checkable sources also record a curve
with thirty exhibited independent rational points, hence rank at least $30$, with attribution to
Claude, Levent Alp\H{o}ge, and Ava Howell appearing on the leaderboard and in subsequent
reports~\cite{DujellaRank2026,MathWorldRank2026}.  The provenance and division of credit are less
formal than the mathematical verification of the displayed objects.

A targeted public search on 23 August 2026 found no paper, preprint, code repository, or
social-media post disclosing a Fable argument for the present \AEconj{} conjecture.  The private work
described in the prompt is therefore not available for comparison.  None of these external AI
claims is used as a mathematical premise below.

\begin{statusbox}{The audit counts and literature dates are factual checks.  The mathematical
content of the report does not depend on any competitive claim.}
\end{statusbox}

\section{Minimal setup and orbit pressure from one counterexample}
\label{sec:setup}

\subsection{The curve form}

Put
\begin{equation}
  \tht:=\log_2 3=\frac{\log 3}{\log 2}=1.584962500721156\ldots .
  \label{eq:theta}
\end{equation}
For $x\in\R$, set
\[
  X=2^x,\qquad Y=3^x.
\]
Then $Y=X^{\tht}$.  The conjecture is exactly the statement
\begin{equation}
  X,Y\in\Z_{>0},\quad Y=X^{\tht}
  \quad\Longrightarrow\quad
  (X,Y)=(2^n,3^n)\text{ for some }n\in\Nzero.
  \label{eq:curve-form}
\end{equation}

The elementary rational case is worth recording because it is used to obtain distinct orbit points.

\begin{lemma}[Rational exponents are integral]
\label{lem:rational-integral}
If $x\in\Q$ and $2^x\in\Z$, then $x\in\Z$.  Consequently every noninteger solution of the two-base
problem is irrational.
\end{lemma}

\begin{proof}
Write $x=a/b$ in lowest terms with $b>0$.  If $2^{a/b}=M\in\Z_{>0}$, then $2^a=M^b$.
Taking the $2$-adic valuation gives $a=bv_2(M)$, so $b\mid a$.  Coprimality forces $b=1$.
\end{proof}

A noninteger solution must also be positive.  Indeed, $2^x$ cannot be a positive integer strictly
between $0$ and $1$.  In fact $2^x\ge2$, hence $x\ge1$.

\subsection{Every counterexample creates linearly many points per dyadic block}

The full primitive-generator theorem from Report III is not needed for the next observation.  A
single irrational solution already supplies the required density.

\begin{theorem}[Orbit pressure in every dyadic block]
\label{thm:block-pressure}
Assume that $x>0$ is a noninteger solution and put $M=2^x\in\Z$, $A=3^x\in\Z$.
For every integer $L\ge0$, the half-open dyadic block
\[
  \mathcal B_L:=\{(X,Y)\in\Z_{>0}^2:2^L\le X<2^{L+1},\ Y=X^{\tht}\}
\]
contains at least
\begin{equation}
  C_L(x):=\ceil{\frac{L+1}{x}}
  \label{eq:block-count}
\end{equation}
distinct points.
\end{theorem}

\begin{proof}
By Lemma~\ref{lem:rational-integral}, $x$ is irrational.  For every integer $k$ satisfying
$0\le k<(L+1)/x$, define
\[
  n_k:=L-\floor{kx}\in\Nzero,
  \qquad s_k:=n_k+kx=L+\{kx\}\in[L,L+1).
\]
Then
\[
  2^{s_k}=2^{n_k}M^k\in\Z,
  \qquad 3^{s_k}=3^{n_k}A^k\in\Z,
\]
so $(2^{s_k},3^{s_k})\in\mathcal B_L$.  If $s_k=s_\ell$ for $k\ne\ell$, then
$(k-\ell)x=n_\ell-n_k\in\Z$, contradicting irrationality.  Thus these points are distinct.
The number of integers $k\ge0$ with $k<(L+1)/x$ is $\ceil{(L+1)/x}$ because $(L+1)/x$ is
not an integer.
\end{proof}

\begin{remark}[Why this passes the control test]
For an integral exponent $x=m$, the same construction collapses: $s_k=L$ for all admissible $k$.
There is only one distinct control point $(2^L,3^L)$ in the block.  Linear block pressure is therefore
a genuinely irrational phenomenon, not a property shared by the known solutions.
\end{remark}

\begin{corollary}[Any successful global exclusion must defeat linear orbit pressure]
\label{cor:linear-pressure}
Any device that assigns a finite ``capacity'' to the integer points it can cover in the $L$th block
must have capacity at least $\asymp_x L$ under a counterexample.  A construction with capacity
$o(L)$ therefore proves the conjecture.
\end{corollary}

The remainder of the report builds an exact, algebraically checkable notion of capacity: the number
of nonzero monomials in a uniformly subunit auxiliary polynomial.

\begin{statusbox}{Theorem~\ref{thm:block-pressure} and its control comparison are \statusproved.
They use only the semigroup closure generated by one solution, not the stronger classification in
Report III.}
\end{statusbox}

\section{Exact dyadic normalization of the near-curve problem}
\label{sec:dyadic}

\subsection{The fixed-function identity}

The general normalization of a block $B\le X<2B$ would replace $X^\tht$ by a scale-dependent
multiple of $z^\tht$.  Choosing $B=2^L$ removes that defect exactly.

\begin{theorem}[Exact dyadic normalization]
\label{thm:dyadic-normalization}
Let $L\in\Nzero$, $u_L=2^L$, $v_L=3^L$, and $f(z)=z^\tht$ on $[1,2]$.  For integers
$X,Y$ with $2^L\le X<2^{L+1}$,
\begin{equation}
  f\!\left(\frac{X}{u_L}\right)-\frac{Y}{v_L}
  =\frac{X^\tht-Y}{3^L}.
  \label{eq:exact-normalization}
\end{equation}
In particular,
\[
  Y=X^\tht
  \quad\Longleftrightarrow\quad
  f(X/u_L)=Y/v_L.
\]
\end{theorem}

\begin{proof}
Since $(2^L)^\tht=3^L$,
\[
  f(X/2^L)=\left(\frac{X}{2^L}\right)^\tht
  =\frac{X^\tht}{3^L}.
\]
Subtract $Y/3^L$.
\end{proof}

This is precisely the input format of Brisebarre and Hanrot~\cite{BrisebarreHanrot2026}: for a fixed
analytic $f:[a,b]\to\R$ and fixed denominators $u,v$, find integers $X,Y$ for which
\begin{equation}
  \left|f(X/u)-Y/v\right|<1/w.
  \label{eq:BH-problem}
\end{equation}
The natural correctly-rounded strip for the present curve is
\begin{equation}
  w_L^{\mathrm{round}}=2\cdot3^L.
  \label{eq:canonical-width}
\end{equation}
Indeed, \eqref{eq:exact-normalization} turns \eqref{eq:BH-problem} into
$|X^\tht-Y|<1/2$, so for each $X$ there is at most one admissible integer $Y$.  Exact points are
contained in every narrower strip as well, so an auxiliary-polynomial algorithm may take $w$ much
larger than $w_L^{\mathrm{round}}$ without losing a genuine solution.

The exponent of the canonical width relative to $u_Lv_L=6^L$ is
\begin{equation}
  \omega:=\frac{\log 3}{\log 6}=0.613147192765458\ldots,
  \qquad 3^L=(6^L)^\omega.
  \label{eq:omega}
\end{equation}

\subsection{A uniform Bernstein ellipse for the power branch}

The fixed function $f(z)=z^\tht$ has a large, explicit complex neighborhood around $[1,2]$.
This is useful both for Chebyshev interpolation and for a directed-rounding certificate.

Let $E_\rho$ be the Bernstein ellipse of parameter $\rho>1$ for $[-1,1]$.  Under the affine map
\[
  z=\frac{3+t}{2},\qquad t\in[-1,1],
\]
the leftmost point of the image ellipse is
\[
  m_\rho=\frac{3-\tfrac12(\rho+\rho^{-1})}{2},
\]
and the maximum modulus is bounded by
\[
  R_\rho=\frac{3+\tfrac12(\rho+\rho^{-1})}{2}.
\]

\begin{lemma}[Uniform analytic domain]
\label{lem:ellipse}
If
\begin{equation}
  1<\rho<3+2\sqrt2,
  \label{eq:rho-range}
\end{equation}
then $m_\rho>0$.  The principal branch of $z^\tht$ is analytic on and inside the corresponding
ellipse around $[1,2]$, and
\begin{equation}
  \max_{z\in E_{\rho,[1,2]}}|z^\tht|\le R_\rho^\tht.
  \label{eq:ellipse-norm}
\end{equation}
\end{lemma}

\begin{proof}
The condition $m_\rho>0$ is equivalent to
$\rho+\rho^{-1}<6$.  The positive root of $\rho^2-6\rho+1=0$ is $3+2\sqrt2$, proving
\eqref{eq:rho-range}.  The whole ellipse then lies in the right half-plane, away from the
branch cut $(-\infty,0]$.  For the principal logarithm,
$|z^\tht|=\exp(\tht\Re\Log z)=|z|^\tht$, and $|z|\le R_\rho$.
\end{proof}

A standard Chebyshev estimate gives the following explicit diagnostic.  If $p_n$ is a degree-$n$
Chebyshev interpolant (or truncation with the same safe constant), then
\begin{equation}
  \norm{f-p_n}_{[1,2],\infty}
  \le \frac{4R_\rho^\tht}{(\rho-1)\rho^n}.
  \label{eq:cheb-error}
\end{equation}
For $\rho=5$, one has $m_5=0.2$, $R_5=2.8$, and therefore
\begin{equation}
  \norm{f-p_n}_\infty\le 2.8^\tht 5^{-n}
  =5.113623057769096\ldots\,5^{-n}.
  \label{eq:rho5-error}
\end{equation}

\begin{table}[ht]
\centering
\small
\begin{tabular}{rcc}
\toprule
Degree $n$ & Safe error bound $E_n$ & Conservative block scale $B$ with $E_n<1/(8B^\tht)$ \\
\midrule
8  & $1.309\times10^{-5}$  & $3.244\times10^{2}$ \\
12 & $2.095\times10^{-8}$  & $1.884\times10^{4}$ \\
20 & $5.362\times10^{-14}$ & $6.354\times10^{7}$ \\
28 & $1.373\times10^{-19}$ & $2.143\times10^{11}$ \\
36 & $3.514\times10^{-25}$ & $7.229\times10^{14}$ \\
48 & $1.439\times10^{-33}$ & $1.416\times10^{20}$ \\
\bottomrule
\end{tabular}
\caption{Uniform one-variable Chebyshev diagnostics at $\rho=5$.  These are not by themselves
Brisebarre--Hanrot lattice-success bounds; they quantify the analytic part of the input.}
\label{tab:cheb}
\end{table}

\subsection{Integer locking by a subunit polynomial}

The key bridge from approximation to arithmetic is elementary but exact.

\begin{theorem}[Subunit integer locking]
\label{thm:integer-locking}
Let $J\subseteq[1,2]$ be an interval, let $P\in\Z[X,Y]$, and let $L\in\Nzero$.
Assume
\begin{equation}
  \sup_{z\in J}\left|P(2^Lz,3^Lz^\tht)\right|<1.
  \label{eq:subunit-graph}
\end{equation}
Then every integer point $(X,Y)$ satisfying
\[
  X/2^L\in J,\qquad Y=X^\tht
\]
is a zero of $P$.
\end{theorem}

\begin{proof}
At such a point, $P(X,Y)\in\Z$.  By \eqref{eq:subunit-graph}, $|P(X,Y)|<1$.  Hence $P(X,Y)=0$.
\end{proof}

The strip version used by Brisebarre--Hanrot is only slightly more general: if
\[
 \sup_{z\in J,\,|t|<1/w}
 \left|P(2^Lz,3^L(z^\tht+t))\right|<1,
\]
then every near point in that strip is a zero.  For the exact conjecture, the $t=0$ slice is enough
for the support-pressure theorem below.

\begin{statusbox}{The normalization, ellipse domain, explicit norm bound, and integer-locking
lemma are \statusproved.  Formula~\eqref{eq:cheb-error} is a standard imported Chebyshev estimate;
the numerical table is \statuscomputed.}
\end{statusbox}

\section{Fewnomial zero budgets and the sparse-certificate criterion}
\label{sec:fewnomial}

\subsection{Exponential monomials form a Chebyshev system}

The next lemma is classical, but its application to the present integer curve appears not to have
been used in the preceding report.

\begin{lemma}[Zero bound for real exponential sums]
\label{lem:exp-zero}
Let $\lambda_1<\cdots<\lambda_T$ be distinct real numbers and let
\[
  h(t)=\sum_{r=1}^T c_r\ee^{\lambda_rt},\qquad c_r\in\R\setminus\{0\}.
\]
Then $h$ has at most $T-1$ real zeros, counted with multiplicity, on every interval.
\end{lemma}

\begin{proof}
Induct on $T$.  The assertion is immediate for $T=1$.  Multiplication by
$\ee^{-\lambda_1t}$ does not change zeros, so consider
\[
  g(t)=c_1+\sum_{r=2}^T c_r\ee^{(\lambda_r-\lambda_1)t}.
\]
If $g$ has $m$ zeros counted with multiplicity on an interval, the multiplicity form of Rolle's
theorem gives at least $m-1$ zeros of $g'$.  The derivative $g'$ is an exponential sum with
$T-1$ distinct exponents, hence has at most $T-2$ zeros by induction.  Thus $m\le T-1$.
\end{proof}

\begin{remark}
The proof is well suited to formalization: it requires no transcendence theorem and no asymptotic
estimate.  It is the standard extended complete Chebyshev property of distinct real exponentials.
\end{remark}

\subsection{A polynomial meets the irrational power curve only finitely often}

For a nonzero polynomial
\[
  P(X,Y)=\sum_{(i,j)\in\cE}c_{ij}X^iY^j,
  \qquad c_{ij}\ne0,
\]
write $T(P)=|\cE|$ for its monomial support size.

\begin{theorem}[Support controls intersections with $Y=X^\tht$]
\label{thm:polynomial-zero-budget}
Let $\tht\notin\Q$ and $P\in\R[X,Y]\setminus\{0\}$.  Then the function
\[
  X\longmapsto P(X,X^\tht),\qquad X>0,
\]
has at most $T(P)-1$ zeros, counted with multiplicity.  In particular, the curve
$Y=X^\tht$ contains at most $T(P)-1$ positive points of the algebraic curve $P(X,Y)=0$.
\end{theorem}

\begin{proof}
Put $X=\ee^t$.  Then
\[
  P(\ee^t,\ee^{\tht t})
  =\sum_{(i,j)\in\cE}c_{ij}\ee^{(i+j\tht)t}.
\]
If $(i,j)\ne(i',j')$ and $i+j\tht=i'+j'\tht$, then
$(j-j')\tht=i'-i$.  Irrationality of $\tht$ forces $j=j'$ and $i=i'$.  The exponents are therefore
distinct, and Lemma~\ref{lem:exp-zero} applies.
\end{proof}

This theorem is stronger than a degree-only Bezout estimate for the present purpose.  It applies to
a single polynomial and measures the exact resource that matters: the number of frequencies
$i+j\tht$.  Total degree $d$ gives only the coarse bound
\begin{equation}
  T(P)\le \Nm_d:=\frac{(d+1)(d+2)}2.
  \label{eq:monomial-count}
\end{equation}

\subsection{Global support pressure under a counterexample}

\begin{theorem}[Global sparse-certificate obstruction]
\label{thm:global-support-pressure}
Assume that $x$ is a noninteger solution.  Let $L\in\Nzero$ and let
$P_L\in\Z[X,Y]\setminus\{0\}$ satisfy
\begin{equation}
  \sup_{1\le z\le2}|P_L(2^Lz,3^Lz^\tht)|<1.
  \label{eq:global-certificate}
\end{equation}
Then
\begin{equation}
  T(P_L)-1\ge \ceil{\frac{L+1}{x}}.
  \label{eq:support-lower}
\end{equation}
Consequently, every such sequence of global certificates has $T(P_L)=\Omega_x(L)$.
\end{theorem}

\begin{proof}
Theorem~\ref{thm:block-pressure} supplies $C_L(x)$ distinct integer points in the $L$th block.
Theorem~\ref{thm:integer-locking} makes all of them zeros of $P_L$.  Theorem
\ref{thm:polynomial-zero-budget} allows at most $T(P_L)-1$ positive zeros.
\end{proof}

\begin{corollary}[Degree threshold]
\label{cor:degree-threshold}
Under the hypotheses of Theorem~\ref{thm:global-support-pressure}, if $P_L$ has total degree at most
$d_L$, then
\begin{equation}
  \frac{(d_L+1)(d_L+2)}2-1
  \ge \ceil{\frac{L+1}{x}}.
  \label{eq:degree-lower}
\end{equation}
In particular, $d_L\ge\sqrt{2L/x}+O_x(1)$.  Therefore the construction of global subunit
polynomials with $d_L=o(\sqrt L)$ would prove the \AEconj{} conjecture.
\end{corollary}

\begin{resultbox}
\textbf{New exact sufficient condition.}
Construct, for an unbounded sequence of block indices $L$, a nonzero $P_L\in\Z[X,Y]$ such that
\[
  \sup_{z\in[1,2]}|P_L(2^Lz,3^Lz^\tht)|<1
\]
and either $T(P_L)=o(L)$ or $\deg P_L=o(\sqrt L)$.  Then no noninteger solution exists.
\end{resultbox}

Unlike many generic analytic properties, this condition does not accidentally eliminate the
integral controls.  A block for the control solutions contains only the endpoint pair
$(2^L,3^L)$, so a two-term polynomial can vanish there without any linear support pressure.

\subsection{The cover-capacity theorem}

A practical auxiliary-polynomial method usually partitions $[1,2]$ into smaller intervals.  The
support criterion survives with an additive capacity.

\begin{definition}[Certified polynomial cover]
For a block index $L$, a certified cover is a finite family
\[
  \{(J_{L,r},P_{L,r}):1\le r\le K_L\}
\]
such that the intervals $J_{L,r}$ cover $[1,2]$, each $P_{L,r}\in\Z[X,Y]\setminus\{0\}$, and
\[
  \sup_{z\in J_{L,r}}|P_{L,r}(2^Lz,3^Lz^\tht)|<1.
\]
Its support capacity is
\begin{equation}
  \operatorname{Cap}_L:=\sum_{r=1}^{K_L}\bigl(T(P_{L,r})-1\bigr).
  \label{eq:capacity}
\end{equation}
\end{definition}

\begin{theorem}[Cover-capacity obstruction]
\label{thm:cover-capacity}
If $x$ is a noninteger solution, every certified cover satisfies
\begin{equation}
  \operatorname{Cap}_L\ge\ceil{\frac{L+1}{x}}.
  \label{eq:capacity-lower}
\end{equation}
If all cover polynomials have total degree at most $d_L$, then
\begin{equation}
  K_L\left(\frac{(d_L+1)(d_L+2)}2-1\right)
  \ge\ceil{\frac{L+1}{x}}.
  \label{eq:Kd2-lower}
\end{equation}
\end{theorem}

\begin{proof}
Assign each of the $C_L(x)$ orbit points from Theorem~\ref{thm:block-pressure} to one cover interval
containing its normalized abscissa.  For a fixed $r$, integer locking makes every assigned point a
zero of $P_{L,r}$, and Theorem~\ref{thm:polynomial-zero-budget} bounds their number by
$T(P_{L,r})-1$.  Sum over $r$ and apply \eqref{eq:monomial-count}.
\end{proof}

\begin{corollary}[Sublinear-capacity criterion]
\label{cor:sublinear-capacity}
The conjecture follows if certified covers exist along an unbounded sequence of block indices with
\[
  \operatorname{Cap}_L=o(L).
\]
A sufficient degree-form condition along that sequence is $K_Ld_L^2=o(L)$.
\end{corollary}

\begin{proof}
Under a noninteger solution, Theorem~\ref{thm:cover-capacity} gives
$\operatorname{Cap}_L/L\ge 1/x+o(1)$ on every block.  This contradicts sublinear capacity along
any unbounded sequence.  The degree-form statement follows from \eqref{eq:Kd2-lower}.
\end{proof}

This is the main quantitative output of the dyadic direction.  It turns a transcendence problem into
a falsifiable approximation-theoretic target.  It also supplies an exact score for experiments:
record not merely the degree and coefficient height of an auxiliary polynomial, but its monomial
support and the total support capacity of the cover.

\subsection{Why support, rather than coefficient height, is the first invariant}

Coefficient height is important for finding a polynomial, but once subunit smallness has been
certified, height disappears from the zero budget.  A polynomial with astronomically large
coefficients but ten surviving monomials still has at most nine positive intersections with the
curve.  Conversely, a dense low-degree polynomial may have enough frequency capacity to absorb the
counterexample orbit.  This suggests three separate optimization axes:

\begin{enumerate}
\item \textbf{analytic size}: make the weighted polynomial uniformly smaller than one;
\item \textbf{support}: keep the number of distinct frequencies $i+j\tht$ below the orbit count;
\item \textbf{cover count}: avoid paying for too many local intervals.
\end{enumerate}

The existing determinant literature usually optimizes the first axis and treats the second as a
byproduct of degree.  Theorem~\ref{thm:cover-capacity} shows that support is itself a proof resource.
A sparse Chebyshev--LLL construction is therefore not merely a speed optimization; it is a logically
stronger proof strategy.

\begin{statusbox}{The exponential zero bound, polynomial intersection theorem, global support
pressure, and cover-capacity criterion are all \statusproved.  The sparse and degree asymptotic
conditions are sufficient conditions, not claimed constructions.}
\end{statusbox}

\section{Conditioning barriers and a concrete lattice pilot}
\label{sec:conditioning}

The support-pressure theorem is conditional: it says what a counterexample would force.  There is
also an unconditional obstruction to constructing the required auxiliary polynomials.  It comes
from the conditioning of a generalized Vandermonde system and is strong enough to locate the
critical degree scale to within a factor of $\sqrt{\log L}$.

\subsection{A support-specific necessary condition}

Let $S\subset\Nzero^2$ be a finite support and put
\[
  \lambda_{ij}:=i+j\tht,\qquad
  \Lambda(S):=\max_{(i,j)\in S}\lambda_{ij},\qquad
  \delta(S):=\min_{\substack{s,s'\in S\\s\ne s'}}|\lambda_s-\lambda_{s'}|.
\]
The separation is positive because $\tht$ is irrational.  For $T=|S|\ge2$, define
\begin{equation}
  \mathfrak B(S):=\frac{T-1}{\log2}
  \log\!\left(
    \frac{1+\exp\!\left(\frac{\log2}{T-1}\Lambda(S)\right)}
         {\frac{\log2}{T-1}\delta(S)}
  \right).
  \label{eq:support-conditioning-budget}
\end{equation}
This is a computable conditioning budget attached only to the frequency support.

\begin{theorem}[Support-conditioning barrier]
\label{thm:support-conditioning}
Let $P\in\Z[X,Y]\setminus\{0\}$ have support $S$.  If
\begin{equation}
  \sup_{1\le z\le2}|P(2^Lz,3^Lz^\tht)|<1,
  \label{eq:conditioning-subunit}
\end{equation}
then $T\ge2$ and
\begin{equation}
  L<\mathfrak B(S).
  \label{eq:conditioning-necessary}
\end{equation}
In particular, no fixed support can furnish global subunit polynomials for arbitrarily large
blocks.
\end{theorem}

\begin{proof}
Write $z=e^t$ and
\[
  H(t):=P(2^Le^t,3^Le^{\tht t})
       =\sum_{s\in S}a_s e^{\lambda_st},
  \qquad a_{ij}:=c_{ij}2^{L\lambda_{ij}},
\]
where the $c_{ij}$ are the nonzero integer coefficients of $P$.  If $T=1$, then $H$ is one
nonzero monomial and its absolute value at $t=0$ is at least one, contrary to
\eqref{eq:conditioning-subunit}.  Thus $T\ge2$.

Take the equally spaced nodes $t_r=rh$, $0\le r<T$, with $h=(\log2)/(T-1)$, and put
$q_s=e^{h\lambda_s}$.  The sampled values satisfy
\[
  (H(t_r))_{0\le r<T}=V(a_s)_{s\in S},
  \qquad V_{r,s}=q_s^r.
\]
This is an ordinary Vandermonde matrix.  The row of $V^{-1}$ recovering $a_s$ is obtained from
\[
  \frac{\prod_{s'\ne s}(U-q_{s'})}{\prod_{s'\ne s}(q_s-q_{s'})}.
\]
The sum of the absolute values of the numerator coefficients is at most
$(1+Q)^{T-1}$, where $Q=\max_sq_s=e^{h\Lambda(S)}$.  Moreover, the mean-value theorem gives
\[
 |q_s-q_{s'}|=|e^{h\lambda_s}-e^{h\lambda_{s'}}|
 \ge h|\lambda_s-\lambda_{s'}|\ge h\delta(S).
\]
Consequently
\begin{equation}
  |a_s|\le
  \left(\frac{1+Q}{h\delta(S)}\right)^{T-1}
  \max_{0\le r<T}|H(t_r)|.
  \label{eq:inverse-vandermonde-bound}
\end{equation}
A subunit polynomial cannot be a nonzero constant, so some nonconstant coefficient occurs.
For that $s$, $\lambda_s\ge1$ and $|a_s|=|c_s|2^{L\lambda_s}\ge2^L$.  Insert this in
\eqref{eq:inverse-vandermonde-bound} and use $\max_r|H(t_r)|<1$ to obtain
\eqref{eq:conditioning-necessary}.
\end{proof}

The theorem identifies the only ways a sparse support can survive far out: its exponential
Vandermonde must become extremely ill-conditioned, which requires either many frequencies, large
frequencies, or very close frequencies.  Merely raising coefficient height cannot evade
\eqref{eq:conditioning-necessary}; coefficient height is already absorbed into the vector $a_s$.

\begin{corollary}[Finite support libraries become extinct]
\label{cor:finite-library}
For every finite collection $\mathscr S$ of supports there is an effective $L_0(\mathscr S)$ such
that no nonzero integer polynomial supported on a member of $\mathscr S$ is globally subunit on a
block $L\ge L_0(\mathscr S)$.
\end{corollary}

\begin{proof}
Take one plus the maximum of $\mathfrak B(S)$ over $S\in\mathscr S$.
\end{proof}

\subsection{An unconditional degree lower bound}

A standard effective linear-form estimate for the two fixed logarithms $\log2$ and $\log3$
gives constants $C_0,C_1>0$ such that
\begin{equation}
  0<|a+b\tht|\quad\Longrightarrow\quad
  |a+b\tht|\ge C_0(\max(2,|a|,|b|))^{-C_1}
  \label{eq:two-log-separation}
\end{equation}
for integers $a,b$; this is an immediate fixed-number specialization of Baker's theory of linear
forms in logarithms~\cite{BakerWustholz,Matveev2000}.  Only the polynomial dependence on the
coefficient bound is used.

\begin{theorem}[Near-critical degree barrier]
\label{thm:unconditional-degree}
There is an effectively computable constant $C>0$ such that the following holds.  If
$P\in\Z[X,Y]\setminus\{0\}$ has total degree at most $d\ge2$ and is globally subunit on block $L$,
then
\begin{equation}
  L\le C d^2\log(d+1).
  \label{eq:degree-conditioning-upper}
\end{equation}
Equivalently, every such polynomial satisfies
\begin{equation}
  d\gg\sqrt{\frac{L}{\log(L+2)}},
  \label{eq:degree-conditioning-lower}
\end{equation}
with an effective absolute implied constant.
\end{theorem}

\begin{proof}
Let $S$ be the support and $T=|S|$.  Then
\[
  T\le\Nm_d=\frac{(d+1)(d+2)}2,
  \qquad \Lambda(S)\le\tht d.
\]
Differences of two frequencies have the form $a+b\tht$ with $|a|,|b|\le d$, so
\eqref{eq:two-log-separation} yields $\delta(S)\ge C_0d^{-C_1}$.  In
\eqref{eq:support-conditioning-budget}, write $h=(\log2)/(T-1)$.  Since
$1+e^{h\Lambda}\le2e^{h\Lambda}$,
\begin{align*}
  (T-1)\log\frac{1+e^{h\Lambda}}{h\delta}
  &\le (T-1)\log2+(T-1)h\Lambda
       +(T-1)\log\frac{T-1}{(\log2)\delta} \\
  &\ll d^2\log(d+1).
\end{align*}
Theorem~\ref{thm:support-conditioning} now gives \eqref{eq:degree-conditioning-upper}.  The
asymptotic inversion is elementary.
\end{proof}

\begin{resultbox}
\textbf{The logarithmic construction window.}
Theorem~\ref{thm:unconditional-degree} says that a global auxiliary polynomial cannot have degree
substantially below $\sqrt{L/\log L}$.  Corollary~\ref{cor:degree-threshold} says that, if a
counterexample exists, every one must have degree at least a constant multiple of $\sqrt L$.
Therefore a construction with
\[
  d_L\asymp\sqrt{\frac{L}{\log L}}
  \quad\text{and}\quad
  \sup_{[1,2]}|P_L(2^Lz,3^Lz^\tht)|<1
\]
would be both close to the unconditional conditioning floor and already sufficient to prove the
conjecture.  The remaining gap is logarithmic rather than qualitative.
\end{resultbox}

\subsection{A reproducible degree-three LLL pilot}
\label{sec:lll-pilot}

To test whether the dyadic normalization is numerically well posed, a small sample lattice was
built from all ten monomials of total degree at most three.  Ten Chebyshev nodes on $[1,2]$ were
used; the evaluation columns were scaled by $10^{30}$, coefficient columns by $10^8$, and SymPy's
exact-integer LLL implementation was applied.  The best of the first ten reduced rows was then
evaluated on a $2001$-point high-precision grid.  The experiment is deliberately modest and is
fully reproduced in Appendix~\ref{app:script}.

\begin{table}[ht]
\centering
\small
\begin{tabular}{rcc}
\toprule
Block $L$ & Smallest sampled supremum & Strictly below one on the sampled grid? \\
\midrule
1 & $6.993\times10^{-5}$ & yes \\
2 & $7.992\times10^{-4}$ & yes \\
3 & $6.302\times10^{-3}$ & yes \\
4 & $5.826\times10^{-2}$ & yes \\
5 & $5.587\times10^{-1}$ & yes \\
6 & $4.099$              & no \\
\bottomrule
\end{tabular}
\caption{Pilot values for a fixed dense degree-three support.  Grid evaluation is not an interval
certificate; the table is computational reconnaissance only.}
\label{tab:lll-pilot}
\end{table}

For example, the row found for $L=5$ was
\begin{align}
P_{5}(X,Y)={}&-18316496+18183588X+9203630X^2+1503225X^3 \notag\\
 &-17089450Y-2644889XY+23548X^2Y \notag\\
 &-676136Y^2-259XY^2+2Y^3.
\label{eq:pilot-polynomial}
\end{align}
Its sampled supremum was $0.5587417310\ldots$.  This is not promoted to a theorem: cancellation
occurs between large terms, so a dense grid alone does not certify the uniform inequality.
Nevertheless, three facts are useful.  First, the scale synchronization is numerically stable
enough for ordinary exact LLL to find strong cancellation.  Second, the rapid deterioration with
$L$ agrees with fixed-support extinction.  Third, it demonstrates why the next implementation
should use the interval-Chebyshev matrix of Brisebarre--Hanrot rather than sample-only evaluation:
their norm test converts an LLL row into a genuine subunit certificate.

\begin{statusbox}{The support-conditioning barrier, fixed-library extinction, and effective
$L\ll d^2\log d$ degree bound are \statusproved, with the standard two-logarithm estimate
\eqref{eq:two-log-separation} imported.  Table~\ref{tab:lll-pilot} and
\eqref{eq:pilot-polynomial} are \statuscomputed and are explicitly not interval certificates.}
\end{statusbox}

\section{Specializing the 2026 near-curve algorithm}
\label{sec:BH}

Brisebarre and Hanrot's June 2026 preprint~\cite{BrisebarreHanrot2026} is the closest recent
algorithmic result to the normalized problem in Section~\ref{sec:dyadic}.  This section separates
three logically different layers:

\begin{enumerate}
\item a rigorous short-vector-to-subunit-polynomial implication;
\item an asymptotic existence and interval-count analysis for dense triangular supports;
\item a heuristic elimination step using two polynomials, whose possible common factor can make the
full point-enumeration algorithm return failure.
\end{enumerate}

For the \AEconj{} conjecture only the first layer is needed.  One certified polynomial per interval
already feeds Theorem~\ref{thm:cover-capacity}; no resultant and no coprimality assumption is
required.

\subsection{The imported short-vector certificate}

For total degree $d$, put
\[
  \Nm_d=\frac{(d+1)(d+2)}2.
\]
Brisebarre--Hanrot form a matrix whose first group of columns records the bivariate Chebyshev
coefficients of
\[
  (z,t)\longmapsto P\bigl(uz,v(f(z)+t)\bigr)
\]
and whose second, diagonal group controls the certified Chebyshev truncation error.  Their
Theorem~4.3 has the following consequence.

\begin{theorem}[Brisebarre--Hanrot short-vector implication, specialized]
\label{thm:BH-short-vector}
Let $J=[a,b]\subset[1,2]$, let $\varepsilon>0$, and take
\[
  f(z)=z^\tht,\qquad u=2^L,\qquad v=3^L.
\]
Choose admissible Bernstein parameters $\rho_1,\rho_2\ge2$ and interpolation orders
$N_1,N_2\ge2$.  If the rounded, tail-augmented Brisebarre--Hanrot matrix $\widehat A$ and an
integer coefficient row $\Lambda=(\lambda_{ij})_{i+j\le d}$ satisfy
\begin{equation}
  \|\Lambda\widehat A\|_2\le
  \frac1{\Nm_d+N_1N_2},
  \label{eq:BH-short-condition}
\end{equation}
then the polynomial
\[
  P(X,Y)=\sum_{i+j\le d}\lambda_{ij}X^iY^j
\]
satisfies
\begin{equation}
  \max_{\substack{z\in J\\ |t|\le\varepsilon}}
  |P(2^Lz,3^L(z^\tht+t))|<1.
  \label{eq:BH-specialized-subunit}
\end{equation}
Consequently every exact integer point $(X,Y)$ with $X/2^L\in J$ and $Y=X^\tht$ is a zero of
$P$.
\end{theorem}

\begin{proof}
The inequality \eqref{eq:BH-specialized-subunit} is exactly Theorem~4.3 of
\cite{BrisebarreHanrot2026} after the substitution from
Theorem~\ref{thm:dyadic-normalization}.  The last assertion follows from integer locking.
\end{proof}

The paper's norm condition is designed to include all rounding and interpolation tails.  Thus a
machine-produced row becomes a mathematical certificate once the matrix enclosure and
\eqref{eq:BH-short-condition} are checked with directed rounding.

\subsection{The proof extends to arbitrary sparse supports}

The triangular monomial set is convenient for degree bounds, but it is not required by the proof of
the short-vector implication.  This is important because the capacity theorem charges actual
support, not ambient degree.

\begin{proposition}[Support-restricted near-curve certificate]
\label{prop:sparse-BH}
Let $S\subset\Nzero^2$ be any finite set of cardinality $T$.  Build the Brisebarre--Hanrot
Chebyshev and tail columns only for the functions
\[
  (z,t)\longmapsto (2^Lz)^i\bigl(3^L(z^\tht+t)\bigr)^j,
  \qquad (i,j)\in S,
\]
with $T\le N_1N_2$, substituting $T$ for $\Nm_d$ throughout the precision and tail
construction.  If the resulting rounded matrix $\widehat A_S$ and
$\Lambda\in\Z^S$ satisfy
\begin{equation}
  \|\Lambda\widehat A_S\|_2\le\frac1{T+N_1N_2},
  \label{eq:sparse-BH-norm}
\end{equation}
then $P_S=\sum_{(i,j)\in S}\lambda_{ij}X^iY^j$ obeys the same uniform subunit conclusion
\eqref{eq:BH-specialized-subunit}.
\end{proposition}

\begin{proof}
In the proof of Theorem~4.3 of~\cite{BrisebarreHanrot2026}, the monomial triangle enters only as an
index set for rows.  The argument bounds the $\ell^1$ norm of the exact Chebyshev coefficient block
plus the diagonal tail block by converting the rounded $\ell^2$ norm.  Deleting rows and replacing
$\Nm_d$ by their number $T$ leaves every inequality unchanged.  Linearity then gives the claimed
polynomial.
\end{proof}

This proposition turns sparse-support search into a rigorous variant of the new algorithm, not an
informal modification.  It also makes the correct objective explicit:
\begin{equation}
  \text{minimize}\quad
  \sum_{r}(T(S_{L,r})-1)
  \quad\text{subject to certified short-vector inequalities.}
  \label{eq:sparse-objective}
\end{equation}
Ordinary LLL quality, coefficient height, and resultant degree are secondary diagnostics.

\subsection{Local ellipse parameters are scale-free}

For $J=[a,b]\subset(0,\infty)$, let $c=(a+b)/2$ and $r=(b-a)/2$.  The Bernstein ellipse with
parameter $\rho$ stays in the right half-plane provided
\[
  c-r\frac{\rho+\rho^{-1}}2>0.
\]
Equivalently, with $S_J=2(a+b)/(b-a)$,
\begin{equation}
  1<\rho<\rho_*(J):=\frac{S_J+\sqrt{S_J^2-4}}2.
  \label{eq:local-rho-star}
\end{equation}
For the full interval $[1,2]$, this gives $\rho_*=3+2\sqrt2$, as in
Lemma~\ref{lem:ellipse}.  On smaller intervals the admissible ellipse grows, but neither
$\rho_*(J)$ nor the bound on $z^\tht$ depends on $L$.  All block growth is isolated in the exact
integer weights $2^{Li}3^{Lj}$.

This separation is computationally valuable.  Analytic enclosures for $z^\tht$ can be cached by
interval shape, while the block index changes only diagonal row scales and required precision.

\subsection{Translation of the dense asymptotic exponents}

The paper introduces a piecewise affine function $g$ satisfying $g(1)=1/2$ and
$g(\lambda)=\frac23\sqrt{2\lambda}+O(1)$.  Its simplified Theorem~4.13 says, heuristically at the
full enumeration layer, that for $\gamma\in[3,\Nm_d]$ a strip of reciprocal width
\[
 w=(uv)^{\frac{d}{3g(\Nm_d/\gamma)}(1+O(1/d))}
\]
can be handled using approximately
\[
 (uv)^{\frac{d}{3g(\Nm_d/\gamma)\gamma}(1+O(1/d))}
\]
subinterval calls, apart from powers of two displayed precisely in the paper.  The interval count
is minimized at the Bombieri--Pila endpoint $\gamma=\Nm_d$.  Since $g(1)=1/2$, define
\begin{equation}
  a_d:=\frac{2d}{3\Nm_d}
      =\frac{4d}{3(d+1)(d+2)}.
  \label{eq:ad}
\end{equation}
With $u=2^L$, $v=3^L$, and $B=2^L$, the dense prediction becomes
\begin{align}
  K_L(d)&=6^{a_dL(1+O(1/d))}
         =B^{\kappa_d(1+O(1/d))},
  \label{eq:BH-call-translation}\\
  \kappa_d&=(1+\tht)a_d,
  \label{eq:kappa-d}\\
  w&=6^{\frac{2d}{3}L(1+O(1/d))}.
  \label{eq:BH-width-translation}
\end{align}
The width in \eqref{eq:BH-width-translation} is much narrower than necessary for exact equality,
but exact points lie in every such strip.

\begin{table}[ht]
\centering
\small
\begin{tabular}{rrrr}
\toprule
$d$ & $\Nm_d$ & $a_d$ relative to $6^L$ & $\kappa_d$ relative to $2^L$ \\
\midrule
4  & 15   & 0.177778 & 0.459549 \\
6  & 28   & 0.142857 & 0.369280 \\
8  & 45   & 0.118519 & 0.306366 \\
10 & 66   & 0.101010 & 0.261107 \\
12 & 91   & 0.087912 & 0.227249 \\
16 & 153  & 0.069717 & 0.180215 \\
20 & 231  & 0.057720 & 0.149204 \\
24 & 325  & 0.049231 & 0.127260 \\
32 & 561  & 0.038027 & 0.098299 \\
40 & 861  & 0.030972 & 0.080061 \\
64 & 2145 & 0.019891 & 0.051418 \\
\bottomrule
\end{tabular}
\caption{Dense-support interval exponents at $\gamma=\Nm_d$.  Values are recomputed in
Appendix~\ref{app:script}.}
\label{tab:BH-exponents}
\end{table}

For fixed $d$, this is an exponential improvement over enumerating all $2^L$ integer abscissae.
If $d=L^\alpha$ with $0<\alpha<1$, the predicted call count is
\begin{equation}
  \exp\!\left(\left(\frac{4\log6}{3}+o(1)\right)L^{1-\alpha}\right),
  \label{eq:BH-subexponential}
\end{equation}
subexponential in the block size $B=2^L$.  The arithmetic input $z^\tht$ is fixed, so there is no
function-growth penalty as $L$ increases.

\subsection{Why the dense algorithm does not yet meet the capacity criterion}

A single call supplies a polynomial with at most $\Nm_d$ monomials.  The dense worst-case support
capacity predicted by \eqref{eq:BH-call-translation} is therefore
\begin{equation}
  \operatorname{Cap}^{\mathrm{dense}}_L(d)
  \asymp \Nm_d\,6^{a_dL}.
  \label{eq:dense-capacity}
\end{equation}
For $d=L^\alpha$ this has logarithm of order $L^{1-\alpha}$ and is far larger than $L$.  Formally
optimizing the model $d^2\exp(cL/d)$ pushes $d$ to order $L$, at which point the support itself is
of order $L^2$.  Thus no choice within the full triangular support makes
\eqref{eq:dense-capacity} sublinear.

This is not a criticism of the near-curve algorithm.  Its objective is to enumerate points faster
than abscissa search, and it succeeds spectacularly at that objective.  The \AEconj{} proof objective is
different: it must beat the linear zero capacity forced by a counterexample.  Proposition
\ref{prop:sparse-BH} shows exactly where to intervene---the monomial support, before lattice
reduction.

\begin{target}[Sparse near-curve frontier]
\label{target:sparse-BH-frontier}
Construct certified covers for the $L$th block from support-restricted matrices such that
\[
  \sum_r(|S_{L,r}|-1)=o(L).
\]
The first concrete benchmark is the global regime
\[
  |S_L|\asymp\frac{L}{\log L},\qquad
  \deg S_L\asymp\sqrt{\frac{L}{\log L}},
\]
which is compatible with the unconditional conditioning floor but contradicts the counterexample
pressure of Theorem~\ref{thm:global-support-pressure}.
\end{target}

The particular support families worth testing are:
\begin{enumerate}
\item anisotropic lower sets adapted to the paper's weighted Taylor order;
\item narrow frequency bands in $i+j\tht$, with line-supported gap powers rejected in advance by
the gap-power obstruction from Report III;
\item supports pruned by the paper's own tail test: its Remark~4.5 forces $\lambda_{ij}=0$ whenever
the certified tail weight of that monomial already exceeds one;
\item unions of a small low-frequency core and a thin high-frequency cancellation shell, scored by
both $\mathfrak B(S)$ and actual LLL output.
\end{enumerate}

\subsection{A proof-oriented certificate format}

A standalone certificate for one interval should contain only exact or outward-rounded data:

\begin{enumerate}
\item $L$ and rational endpoints $a,b$;
\item the support $S$ and integer coefficient vector $\Lambda$;
\item rational enclosures for $\tht$, ellipse maxima, and every required Chebyshev coefficient;
\item the rounded integer matrix $\widehat A_S$ and the precision exponent used to create it;
\item a rational upper bound for $\|\Lambda\widehat A_S\|_2^2$ strictly below
$(T+N_1N_2)^{-2}$;
\item the analytic tail inequality used in constructing the diagonal columns.
\end{enumerate}

The checker does not run LLL.  It reconstructs the enclosures, verifies the norm inequality, and
then invokes Proposition~\ref{prop:sparse-BH}.  A collection of such records is accepted only after
checking that its rational intervals cover $[1,2]$ and summing the exact support capacity.
This division of labor is well suited to Lean: untrusted high-performance code proposes rows;
a small kernel checks the certificate.

\begin{statusbox}{The exact normalization and support-restricted short-vector implication are
\statusproved, using the imported Brisebarre--Hanrot norm theorem.  The exponent translation and
tables are \statuscomputed from their asymptotic formulas.  The full enumeration estimate is
heuristic only because their two auxiliary polynomials may share a factor; the one-polynomial
certificate route used here has no coprimality heuristic.  The sublinear sparse frontier is a
\statustarget.}
\end{statusbox}

\section{The missing operator: a toric no-go and a Carlitz model}
\label{sec:operator}

The ordinary logarithmic form of the problem contains a rank-one determinant:
\[
  \log2\,\log A-\log3\,\log M=0.
\]
The challenge is not to discover this relation but to prove that it comes from linear relations.
Positive characteristic has an exact theorem of this kind.  Before developing it, it is useful to
show why the most obvious characteristic-zero scaling operator has no room to create such a
lifting principle.

\subsection{Diagonal toric scaling cannot mix monomials}

Let $\cT$ be the pullback operator on $\C[X,Y]$ defined by
\[
  (\cT P)(X,Y):=P(2X,3Y).
\]

\begin{theorem}[Polynomial toric semi-invariants]
\label{thm:toric-semi-invariant}
If $P\in\C[X,Y]\setminus\{0\}$ and $c\in\C$ satisfy
\[
  P(2X,3Y)=cP(X,Y),
\]
then $P$ is a single monomial.  More generally, every finite-dimensional $\cT$-invariant subspace
of $\C[X,Y]$ is spanned by monomials.
\end{theorem}

\begin{proof}
Write $P=\sum c_{ij}X^iY^j$.  Comparing coefficients gives
$c=2^i3^j$ for every $(i,j)$ in the support.  Unique factorization makes the eigenvalues
$2^i3^j$ pairwise distinct, so the support has one element.

For the second statement, take $P$ in a finite-dimensional invariant subspace $V$.  Its finitely
many monomial components are eigenvectors of $\cT$ with distinct eigenvalues.  Lagrange
interpolation in the operator $\cT|_V$ projects $P$ onto each of these eigenspaces, so every
monomial component of $P$ belongs to $V$.  Applying this to a basis of $V$ proves the claim.
\end{proof}

Thus the map $(X,Y)\mapsto(2X,3Y)$ is semisimple in the least helpful possible way.  It preserves
finite-dimensional polynomial modules only by separating monomials; it never creates a nontrivial
coupled system.  The toric-mask analysis in Report III reaches a related conclusion for exponent
lattices.  The point here is operator-theoretic: a successful relation-lifting theorem must enlarge
the function space and introduce a semilinear or inhomogeneous action, rather than merely iterate
geometric scaling.

\subsection{A linear-algebra lemma for determinant relations}

The next elementary classification is the bridge between algebraic-independence theorems and
multiplier rigidity.

\begin{lemma}[A rational subspace of singular $2\times2$ matrices]
\label{lem:singular-subspace}
Let $K$ be a field and let $W\subset M_2(K)$ be a $K$-linear subspace all of whose matrices are
singular.  Suppose the map taking a matrix to its first row is onto $K^2$.  Then there is a unique
$q\in K$ such that
\begin{equation}
  W=\left\{
  \begin{pmatrix}s&t\\qs&qt\end{pmatrix}:s,t\in K
  \right\}.
  \label{eq:compression-space}
\end{equation}
\end{lemma}

\begin{proof}
Choose $A,B\in W$ with first rows $(1,0)$ and $(0,1)$.  Singularity gives
\[
 A=\begin{pmatrix}1&0\\a&0\end{pmatrix},
 \qquad
 B=\begin{pmatrix}0&1\\0&d\end{pmatrix}.
\]
Every $sA+tB$ is singular, so $st(d-a)=0$ for all $s,t$ and $d=a=:q$.
If $C\in W$ has zero first row, write its second row as $(u,v)$.  Singularity of $A+rC$ for all
$r$ gives $v=0$, and singularity of $B+rC$ gives $u=0$.  Hence the first-row map is injective as
well as surjective, and \eqref{eq:compression-space} follows.
\end{proof}

Let $z=(z_1,z_2,z_3,z_4)$ be a tuple in a field extension of $K$.  Denote by $I_1(z)$ the space of
$K$-linear forms vanishing at $z$, and by $\langle I_1(z)\rangle$ the ideal they generate.

\begin{proposition}[Quadratic descent implies common scalar]
\label{prop:quadratic-descent}
Assume $z_1,z_2$ are $K$-linearly independent and
\[
  z_1z_4-z_2z_3=0.
\]
If there is a field extension $K'/K$ such that the determinant polynomial
\begin{equation}
  X_1X_4-X_2X_3
  \label{eq:determinant-polynomial}
\end{equation}
belongs to $\langle I_1(z)\rangle K'[X_1,X_2,X_3,X_4]$, then there is $q\in K$ with
\[
  z_3=qz_1,
  \qquad z_4=qz_2.
\]
\end{proposition}

\begin{proof}
Choose a $K$-basis of the span of the four $z_i$ and let $W\subset K^4\cong M_2(K)$ be the
smallest $K$-linear subspace defined by their linear relations.  Membership of
\eqref{eq:determinant-polynomial} in the scalar-extended linear-relation ideal says that the
determinant vanishes identically on $W\otimes_K K'$, and hence on $W$.
Independence of $z_1,z_2$ makes the first-row projection of $W$ onto $K^2$.
Lemma~\ref{lem:singular-subspace} gives the stated form.
\end{proof}

\subsection{Carlitz logarithms and their complete relation ideal}

Let
\[
  k=\Fq(\vartheta),\qquad A=\Fq[\vartheta],
\]
and let $\C_\infty$ be the completion of an algebraic closure of the completion of $k$ at infinity.
The Carlitz exponential is
\[
  \exp_C(z)=\sum_{n\ge0}\frac{z^{q^n}}{D_n},
\]
with its standard Carlitz denominators $D_n$.  Put
\begin{equation}
  \cL_C:=\{\lambda\in\C_\infty:\exp_C(\lambda)\in\overline{k}\}.
  \label{eq:Carlitz-log-set}
\end{equation}
Papanikolas proved that a $k$-linearly independent finite family in $\cL_C$ is algebraically
independent over $\overline{k}$~\cite{Papanikolas2008}.

That theorem immediately determines \emph{all} polynomial relations in an arbitrary finite
family, not only independent ones.

\begin{theorem}[Relation ideal for algebraic Carlitz logarithms]
\label{thm:Carlitz-relation-ideal}
Let $\lambda_1,\ldots,\lambda_n\in\cL_C$.  The ideal of polynomial relations over $\overline{k}$
among the $\lambda_i$ is generated by their $k$-linear relations.
\end{theorem}

\begin{proof}
Choose a maximal $k$-linearly independent subfamily
$\mu_1,\ldots,\mu_r$.  Every $\lambda_i$ is a $k$-linear combination of the $\mu_j$, so the
corresponding linear forms generate an ideal $J$ contained in the relation ideal.  Reduction modulo
$J$ turns any polynomial in the $\lambda_i$ into a polynomial in the $\mu_j$.  Papanikolas's
theorem says the $\mu_j$ are algebraically independent over $\overline{k}$, so a reduced
polynomial can vanish only if it is zero.  Hence the full relation ideal is $J$.
\end{proof}

\begin{theorem}[Two-direction Carlitz multiplier rigidity]
\label{thm:Carlitz-multiplier}
Let $\lambda_1,\lambda_2\in\cL_C$ be $k$-linearly independent.  Let $c\in\C_\infty$ and suppose
$c\lambda_1,c\lambda_2\in\cL_C$.  Then
\begin{equation}
  c\in k.
  \label{eq:Carlitz-multiplier-rational}
\end{equation}
\end{theorem}

\begin{proof}
Apply Theorem~\ref{thm:Carlitz-relation-ideal} to the four Carlitz logarithms
\[
  (\lambda_1,\lambda_2,c\lambda_1,c\lambda_2).
\]
Their determinant relation
\[
  \lambda_1(c\lambda_2)-\lambda_2(c\lambda_1)=0
\]
belongs to the ideal generated by $k$-linear relations.  Proposition
\ref{prop:quadratic-descent}, with $K=k$, gives $q\in k$ such that
$c\lambda_i=q\lambda_i$ for $i=1,2$.  Since the $\lambda_i$ are nonzero, $c=q$.
\end{proof}

\begin{remark}[Exact positive-characteristic analogue]
The theorem has precisely the desired two-direction form.  Two independent logarithmic directions
are given; multiplication by the same scalar preserves the class of logarithms of algebraic
objects; the scalar is forced into the rational function field.  No height estimate, density
argument, or genericity hypothesis is involved.
\end{remark}

\subsection{The exact characteristic-zero quadratic target}

Suppose $x$ is a solution and set
\[
  \ell_1=\log2,\qquad \ell_2=\log3,
  \qquad \ell_3=\log M=x\log2,
  \qquad \ell_4=\log A=x\log3.
\]
The first two logarithms are $\Q$-linearly independent, and
\begin{equation}
  \ell_1\ell_4-\ell_2\ell_3=0.
  \label{eq:ordinary-log-determinant}
\end{equation}
Let $I_1(\ell)$ denote the rational linear relations among these four numbers.

\begin{target}[Quadratic logarithm relation target]
\label{target:quadratic-log}
Prove, for the above four logarithms, that
\begin{equation}
  X_1X_4-X_2X_3\in
  \big\langle I_1(\ell)\big\rangle\overline{\Q}[X_1,X_2,X_3,X_4].
  \label{eq:quadratic-target}
\end{equation}
\end{target}

\begin{theorem}[The quadratic target settles the conjecture]
\label{thm:quadratic-target-closes}
If Target~\ref{target:quadratic-log} holds for every putative solution, then the \AEconj{} conjecture
is true.
\end{theorem}

\begin{proof}
Let $W$ be the $\Q$-linear subspace defined by the rational linear relations among the $\ell_i$.
Condition \eqref{eq:quadratic-target} makes the determinant vanish identically on
$W\otimes_\Q\overline\Q$.  Proposition~\ref{prop:quadratic-descent}, over $K=\Q$, gives
$q\in\Q$ with
\[
  \log M=q\log2,\qquad \log A=q\log3.
\]
Thus $x=q\in\Q$, and Lemma~\ref{lem:rational-integral} gives $x\in\Z$.
\end{proof}

This target is much weaker than Schanuel's conjecture.  It asks about one explicitly known
quadratic relation in one four-element family; it does not ask for a transcendence-degree formula
for arbitrary exponentials.  Schanuel would imply it: a $\Q$-basis of logarithms of algebraic
numbers would be algebraically independent, so all polynomial relations would be generated by the
linear ones.  Papanikolas proves exactly that relation-ideal property for Carlitz logarithms.

The target also clarifies what a partial theorem must accomplish.  A lower bound for a nonzero
quadratic expression is irrelevant because the determinant is identically zero.  One must instead
show that a vanishing quadratic tensor has linear provenance.

\subsection{What Estienne's 2026 Mahler proof adds}

Papanikolas's original proof uses Tannakian Galois groups of $t$-motives.  Estienne's March 2026
paper~\cite{Estienne2026} supplies a second mechanism: interpolate the Carlitz logarithms by
functions $G_i(Z)$ satisfying a finite Mahler system and transfer algebraic independence from the
functions to an algebraic specialization.  In simplified form, the interpolants satisfy an
inhomogeneous relation
\begin{equation}
  G_i(Z^{q^d})=\pi_d(Z)\bigl(G_i(Z)-M_i(Z)\bigr),
  \label{eq:Estienne-functional}
\end{equation}
which becomes a homogeneous finite-dimensional system after adjoining the necessary auxiliary
functions, including an interpolant for the Carlitz period.  A positive-characteristic analogue of
Nishioka's specialization theorem then preserves transcendence degree.

This proof exposes four pieces that a characteristic-zero construction would need:

\begin{enumerate}
\item a finite vector of analytic functions whose algebraic specialization contains
$\log2,\log3,\log M,\log A$;
\item a nontrivial endomorphism, analogous to $Z\mapsto Z^{q^d}$, under which that vector satisfies
a rational or algebraic matrix equation;
\item a functional relation theorem saying that the determinant relation is generated by linear
functional relations;
\item a regular-specialization theorem that preserves this relation ideal at the desired algebraic
point.
\end{enumerate}

The toric operator of Theorem~\ref{thm:toric-semi-invariant} fails the second item: it diagonalizes
rather than couples.  The next calculation shows why the most obvious ordinary-log Mahler family
also fails finite closure.

\subsection{The root-alphabet obstruction for the naive ordinary logarithm family}

For an algebraic parameter $c$, define near the origin
\[
  F_c(z):=-\log(1-cz)=\sum_{n\ge1}\frac{c^nz^n}{n}.
\]
At $z=1$ this contains the relevant logarithms:
\[
  F_{1/2}(1)=\log2,
  \qquad F_{2/3}(1)=\log3,
  \qquad F_{1-1/M}(1)=\log M.
\]
For an integer $q\ge2$, choose $\alpha$ with $\alpha^q=c$.  The local factorization
\[
  1-cz^q=\prod_{\zeta^q=1}(1-\alpha\zeta z)
\]
gives the exact germ identity
\begin{equation}
  F_c(z^q)=\sum_{\zeta^q=1}F_{\alpha\zeta}(z).
  \label{eq:root-alphabet}
\end{equation}
Repeated pullback therefore introduces the parameters
\[
  c^{1/q^n}\zeta,
  \qquad \zeta^{q^n}=1,
\]
for arbitrarily large $n$.  Starting from $c=1/2$ or $2/3$, these are infinitely many distinct
algebraic parameters.  Hence the finite list
\[
  F_{1/2},F_{2/3},F_{1-1/M},F_{1-1/A}
\]
is not contained in any evident finite-dimensional module under $z\mapsto z^q$.

\begin{warning}
Equation~\eqref{eq:root-alphabet} rules out only the naive parameter-fixed Mahler module.  It does
not prove that no characteristic-zero operator exists.  A successful construction may use several
variables, parameter dynamics, differential--difference coupling, a noncommutative module, or a
special function other than $-\log(1-cz)$.
\end{warning}

\subsection{A concrete operator-lifting research target}

The function-field theorem suggests the following deliberately narrow program.

\begin{target}[Finite operator module for four logarithms]
\label{target:operator-module}
Construct an algebraic dynamical system $\sigma:V\dashrightarrow V$, a regular algebraic point
$\alpha\in V(\overline\Q)$, and a finite vector of analytic functions $\mathbf F$ such that:
\begin{enumerate}
\item $\mathbf F(\sigma z)=A(z)\mathbf F(z)$ for
$A(z)\in\mathrm{GL}_n(\overline\Q(V))$ after adjoining finitely many inhomogeneous coordinates;
\item the four ordinary logarithms occur as $\overline\Q$-linear combinations of
$\mathbf F(\alpha)$;
\item every degree-two functional relation among the corresponding four components is generated by
functional linear relations;
\item specialization at $\alpha$ preserves that degree-two relation statement.
\end{enumerate}
Only the determinant relation \eqref{eq:ordinary-log-determinant} needs to be lifted.
\end{target}

A useful relaxation is to allow an infinite ambient family but prove that the determinant lies in a
finite invariant quotient.  The root-alphabet identity then becomes input rather than an immediate
obstruction: one would seek a finite quotient of the rooted parameter tree that retains the four
special values and their quadratic tensor.

\begin{statusbox}{The toric semi-invariant theorem, singular-subspace classification, quadratic
descent proposition, Carlitz relation-ideal theorem, and Carlitz multiplier-rigidity theorem are
\statusproved.  Papanikolas's algebraic-independence theorem and Estienne's Mahler construction are
\statusimported.  The ordinary quadratic relation and finite operator module are
\statustarget.  The root-alphabet calculation is an exact \statusnegative audit of one naive
module, not a general impossibility theorem.}
\end{statusbox}

\section{A differential logarithm lift, and why it specializes incorrectly}
\label{sec:differential-audit}

The root-alphabet obstruction in Section~\ref{sec:operator} concerns a difference
operator.  Ordinary logarithms also admit an elementary differential lift.  At first sight it
looks even better: the functions satisfy a finite rational differential system, and their
functional algebraic relations can be determined completely.  The difficulty is that this lift
has \emph{too few} functional relations.  Multiplicative identities among the special values are
created only at the specialization point.

\subsection{Fixed-endpoint logarithm germs}

For a rational number $q>1$, put
\begin{equation}
  c_q:=1-\frac1q,
  \qquad
  \cL_q(z):=-\log(1-c_qz),
  \label{eq:differential-log-family}
\end{equation}
where the branch is the germ at $z=0$.  Then
\begin{equation}
  \cL_q(1)=\log q,
  \qquad
  \cL_q'(z)=\frac{c_q}{1-c_qz}.
  \label{eq:differential-log-values}
\end{equation}
Thus $(1,\cL_{q_1},\ldots,\cL_{q_n})^t$ satisfies an inhomogeneous unipotent differential
system with coefficients in $\Q(z)$, made homogeneous by adjoining the constant component.
The point $z=1$ is ordinary because $1/c_q=q/(q-1)>1$.

The complete functional relation theorem is especially simple.

\begin{theorem}[Monodromy independence of the fixed-endpoint logarithms]
\label{thm:monodromy-independence}
Let $q_1,\ldots,q_n>1$ be distinct rational numbers.  Then the germs
\[
  \cL_{q_1}(z),\ldots,\cL_{q_n}(z)
\]
are algebraically independent over $\C(z)$.
\end{theorem}

\begin{proof}
The singular points $a_i=1/c_{q_i}$ are distinct.  Suppose that there is a nonzero polynomial
relation
\[
  P\bigl(z,\cL_{q_1}(z),\ldots,\cL_{q_n}(z)\bigr)=0,
  \qquad P\in\C(z)[Y_1,\ldots,Y_n],
\]
and choose one of minimal total degree.  Analytic continuation around a small loop enclosing
$a_i$ and no other singularity replaces $\cL_{q_i}$ by $\cL_{q_i}+\omega$, where
$\omega=\pm2\pi i$, and leaves all other germs unchanged.  Hence
\[
  \Delta_iP:=P(z,Y_1,\ldots,Y_i+\omega,\ldots,Y_n)-P(z,Y_1,\ldots,Y_n)
\]
vanishes after substitution of the germs.  If $P$ depends on $Y_i$, then $\Delta_iP$ is a
nonzero polynomial relation of smaller total degree, a contradiction.  Thus $P$ is independent
of $Y_i$.  Repeating this for every $i$ makes $P$ an element of $\C(z)$, which cannot vanish
identically after substitution unless $P=0$.
\end{proof}

If parameters are repeated, the full relation ideal is obtained by adjoining the obvious linear
relations $Y_i-Y_j$ for equal $q_i=q_j$ and then applying the theorem to one representative of
each distinct parameter.

\subsection{The specialization defect already appears at integral controls}

Theorem~\ref{thm:monodromy-independence} does not transfer to $z=1$.  For example,
\begin{equation}
  \cL_4(1)-2\cL_2(1)=0,
  \qquad
  \cL_9(1)-2\cL_3(1)=0,
  \label{eq:differential-integral-collapse}
\end{equation}
while the three distinct germs $\cL_2,\cL_4,\cL_9$ are algebraically independent over
$\C(z)$.  More generally, every multiplicative relation
$\prod_iq_i^{m_i}=1$ gives the value relation
\begin{equation}
  \sum_i m_i\cL_{q_i}(1)=0,
  \label{eq:value-multiplicative-relation}
\end{equation}
although no corresponding functional relation exists when the $q_i$ are distinct.

For a candidate pair $(M,A)$ define
\begin{equation}
  D_{M,A}(z):=\cL_M(z)\cL_3(z)-\cL_A(z)\cL_2(z).
  \label{eq:differential-determinant}
\end{equation}
If $2^x=M$ and $3^x=A$, then $D_{M,A}(1)=0$.  On the other hand,
$D_{M,A}(z)$ is not the zero germ for any nonintegral candidate.  Indeed, if
$M\notin\{2,3,A\}$ this follows from algebraic independence of the four germs.  If $M=3$,
then $D_{M,A}=\cL_3^2-\cL_A\cL_2$ is nonzero by algebraic independence of the three distinct
germs.  The output $A$ is larger than both $3$ and $M$ for a nonintegral candidate, so there
are no other coincidences.

Even the valid control $x=2$, $(M,A)=(4,9)$, shows the correct mechanism:
$D_{4,9}(1)=0$ because the two linear relations in
\eqref{eq:differential-integral-collapse} generate the determinant relation at the value level,
but neither linear relation lifts to the function level.  Therefore no specialization theorem
that simply preserves the functional relation ideal of the family
\eqref{eq:differential-log-family} can prove the desired quadratic descent: it would reject the
integral controls as well.

\begin{proposition}[Exact failure mode of the naive differential lift]
\label{prop:differential-failure}
The family $\{\cL_q\}$ has a finite rational differential system and a completely rigid
functional relation ideal, but evaluation at $z=1$ is not relation-faithful even in degree one.
Consequently it cannot be the function family in Target~\ref{target:operator-module} without
adding structure that lifts multiplication of the parameters before specialization.
\end{proposition}

\begin{proof}
The differential system is \eqref{eq:differential-log-values}; functional rigidity is
Theorem~\ref{thm:monodromy-independence}; failure in degree one is
\eqref{eq:differential-integral-collapse}.
\end{proof}

The missing structure is now clearer.  A useful lift must remember that
$\log(qr)=\log q+\log r$ at the \emph{functional or motivic} level, not merely after evaluation.
Kummer $1$-motives do encode this functoriality, but their unconditional period theory controls
linear relations, whereas the determinant lies in a tensor square.  The Carlitz setting succeeds
because the Frobenius difference module controls both the linear relation space and the full
polynomial relation ideal.

\begin{target}[Relation-faithful logarithm lift]
\label{target:relation-faithful-lift}
For the finitely generated subgroup of $\Q_{>0}^{\times}$ spanned by the prime divisors of
$6MA$, construct analytic objects $\mathscr K_q$ with algebraic operator data and an algebraic
specialization $\alpha$ such that:
\begin{enumerate}
\item $\mathscr K_q(\alpha)=\log q$;
\item every multiplicative relation among the $q$ lifts to the corresponding functional linear
relation among the $\mathscr K_q$;
\item degree-two functional relations are generated by those lifted linear relations; and
\item specialization at $\alpha$ introduces no new degree-two relation.
\end{enumerate}
For the conjecture only the determinant tensor is needed, not full algebraic independence in all
degrees.
\end{target}

\begin{statusbox}{The monodromy independence theorem and the specialization-defect proposition
are \statusproved.  They give a second exact no-go, independent of the Mahler root-alphabet
obstruction.  Target~\ref{target:relation-faithful-lift} is \statustarget.}
\end{statusbox}

\section{Rank-two logarithmic grids and the real additive threshold}
\label{sec:harrison}

Report III already established the carry-language lower bounds.  The purpose of this section is
not to repeat their construction.  It is to place them against the precise July 2026 theorem of
Harrison~\cite{Harrison2026} and determine which kind of strengthening could actually contradict
a counterexample.

\subsection{The triangular smooth grid}

For $R\ge0$ define
\begin{equation}
  \cG_R:=\{2^a3^b:a,b\in\Nzero,\ a+b\le R\}.
  \label{eq:smooth-grid}
\end{equation}
Unique factorization gives
\begin{equation}
  |\cG_R|=\frac{(R+1)(R+2)}2,
  \qquad
  \cG_R\subseteq\{1,2,\ldots,3^R\}.
  \label{eq:smooth-grid-size}
\end{equation}
Its logarithm is the triangular rank-two generalized arithmetic progression
\[
  \log\cG_R=\{a\log2+b\log3:a,b\ge0,\ a+b\le R\}.
\]
Thus $\cG_R$ is polynomially large in $R$ but only of size $\asymp(\log N)^2$ inside the
ambient $N$-term progression with $N=3^R$.  It lies exactly at the intersection of two frontiers
identified in Harrison's paper: sparse subsets and explicit rank or dimension dependence.

If $x$ were a counterexample and $M=2^x$, $A=3^x$, then
\begin{equation}
  \cG_R^{[x]}=\{M^aA^b:a,b\ge0,\ a+b\le R\}\subset\Z.
  \label{eq:powered-smooth-grid}
\end{equation}
The same triangular exponent set survives, but the two geometric ratios become integers.

\subsection{What Harrison's theorem gives on this grid}

For fixed integers $1\le r\le s$, put
\[
  \eta(r,s):=\max\{1,\min(r,s-1)\}.
\]
Harrison's Theorem~1.3 states that, for irrational $c$, the number of ordered non-permutation
solutions of
\[
  \sum_{i=1}^{s}u_i^c=\sum_{j=1}^{r}v_j^c
\]
with all variables in a subset of an $N$-term arithmetic progression is
$O_s(|B|^{\eta(r,s)}(\log N)^{C_2(s)})$, where $C_2(s)$ is effective.
The immediate specialization is the following.

\begin{corollary}[Fixed-arity Harrison bound on the smooth grid]
\label{cor:harrison-grid}
Let $c$ be irrational and fix $1\le r\le s$.  The number of ordered non-permutation solutions
with all variables in $\cG_R$ is
\begin{equation}
  O_s\!\left(R^{2\eta(r,s)+C_2(s)}\right)
  \label{eq:harrison-grid-bound}
\end{equation}
as $R\to\infty$.
\end{corollary}

\begin{proof}
Apply Harrison's theorem with $B=\cG_R$ and the ambient progression
$\{1,\ldots,3^R\}$.  Then $|B|\asymp R^2$ and $\log N=R\log3$.
\end{proof}

This is a genuine restriction for every fixed arity.  It is nevertheless compatible with all
counterexample carries.  A particularly transparent calibration comes from the two primitive
moves.  The two orders of expanding $(6t)^x$ give, for every $t\in\cG_{R-1}$,
\begin{equation}
  (M-1)(3t)^x+A\,t^x=(A-1)(2t)^x+M\,t^x.
  \label{eq:plaquette-additive-relation}
\end{equation}
Both sides have $M+A-1$ terms, and the relation is non-permutational.  Hence there are
$\asymp R^2$ scaled instances at this fixed arity.  Corollary~\ref{cor:harrison-grid} allows a
much larger polynomial, so there is no tension.

\subsection{From total energy to a fixed-target fiber}

For a real exponent $c$, an integer $n\ge1$, and a real target $T$, let
\begin{equation}
  \cR_{n,R}^{(c)}(T):=
  \#\left\{(m_g)_{g\in\cG_R}\in\Nzero^{\cG_R}:
  \sum_gm_g=n,\quad \sum_gm_gg^c=T\right\}.
  \label{eq:unordered-fiber}
\end{equation}
This counts unordered multiset representations.  Harrison's ordered energy bound yields a
fixed-target consequence, but its shape is too large in the arity needed here.

\begin{proposition}[Square-root fiber consequence]
\label{prop:harrison-fiber}
Fix an integer $n\ge2$ and an irrational $c$.  Then
\begin{equation}
  \max_T\cR_{n,R}^{(c)}(T)
  \ll_n R^{n-1+C_2(n)/2}
  \label{eq:harrison-fiber-bound}
\end{equation}
for $R\ge2$.
\end{proposition}

\begin{proof}
Choose one canonical ordering of each multiset in a target fiber.  If the fiber contains $H$
distinct multisets, then ordered pairs of distinct canonical representatives give
$H(H-1)$ non-permutation solutions of the equal-arity equation.  Harrison's theorem with
$r=s=n$ bounds their number by
\[
  O_n\bigl(|\cG_R|^{n-1}(\log3^R)^{C_2(n)}\bigr).
\]
Use $|\cG_R|\asymp R^2$ and take square roots.
\end{proof}

The dependence on fixed $n$ is essential.  More importantly, even a hypothetical uniform version
of the same \emph{shape} would not be enough.  If $n\asymp R$, the factor
$R^{n-1}$ in \eqref{eq:harrison-fiber-bound} is
\begin{equation}
  \exp\bigl((1+o(1))n\log R\bigr),
  \label{eq:energy-shape-wall}
\end{equation}
whereas the inherited carry fibers grow only exponentially in $n$.  Making $C_2(n)$ and the
implicit constant explicit cannot remove the $n\log R$ term.  The missing estimate must exploit
the fixed target and the boundary geometry, not merely count all equal sums.

\subsection{The inherited exponential boundary fiber}

For a hypothetical counterexample, Report III defines $\tau_x>0$ by
\begin{equation}
  \sum_{(a,b)\in\Nzero^2\setminus\{(0,0)\}}
  \exp\bigl(-\tau_x(M^aA^b-1)\bigr)=1
  \label{eq:imported-carry-pressure}
\end{equation}
and proves that, along the appropriate residue class of arities, there are targets
$T_n=M^{r_n}A^{s_n}$ with $r_n+s_n\le n$ such that
\begin{equation}
  \cR_{n+1,n}^{(x)}(T_n)
  \gg_x \frac{\exp(\tau_xn)}{n^2}.
  \label{eq:imported-boundary-fiber}
\end{equation}
This statement is imported rather than reproved.

There is also a simpler mixed-direction calibration.  Take $r=s=m$ primitive unit moves.  The
$\binom{2m}{m}$ possible orders give distinct multiset representations of the same target
$M^mA^m$, all using
\begin{equation}
  n_m=1+m(M+A-2)
  \label{eq:balanced-carry-arity}
\end{equation}
terms from $\cG_{2m}^{[x]}$.  Therefore
\begin{equation}
  \liminf_{m\to\infty}\frac1{n_m}
  \log\cR_{n_m,2m}^{(x)}(M^mA^m)
  \ge \frac{2\log2}{M+A-2}.
  \label{eq:elementary-mixed-entropy}
\end{equation}
This lower rate isolates the genuinely two-directional source of entropy: one primitive move in
each direction has only one order, while interleaving the two creates the central binomial
coefficient.

\begin{target}[Boundary-conditioned mixed-fiber theorem]
\label{target:boundary-fiber}
Develop a rank-sensitive theorem for the triangular log-grid that bounds fixed-target fibers of
arity $n\asymp R$ after isolating their character-balanced algebraic parts.  It must do at least
one of the following:
\begin{enumerate}
\item give, for the smooth boundary targets $(2^r3^s)^c$, an upper exponential rate strictly
below the mixed rate forced by \eqref{eq:elementary-mixed-entropy}; or
\item classify all positive-dimensional component signatures and prove that two independent
primitive carry directions cannot coexist for irrational $c$ with integer coefficients.
\end{enumerate}
A uniform extension of \eqref{eq:harrison-grid-bound} with the same power of $|\cG_R|$ is not
sufficient.
\end{target}

Harrison's proof explains where such a theorem must differ.  Transcendental rational points are
controlled by o-minimal counting, while semi-algebraic curves are treated through functional
transcendence and induction.  The counterexample carries are character-balanced and therefore
live in the algebraic part after a common dilation parameter is introduced.  Their entropy can be
hidden in the number of algebraic component signatures as arity grows.  The required advance is
thus a uniform classification or compression of those signatures, not a stronger estimate for
the transcendental part alone.

\begin{statusbox}{Corollary~\ref{cor:harrison-grid} and
Proposition~\ref{prop:harrison-fiber} are \statusproved from Harrison's imported theorem.
The pressure lower bound is imported from Report III, and
\eqref{eq:elementary-mixed-entropy} is its elementary two-direction specialization.
Target~\ref{target:boundary-fiber} is \statustarget.}
\end{statusbox}

\section{Route triage and integral-control failure tests}
\label{sec:audits}

A proposed obstruction is useful only if it survives every integral exponent.  This section
collects the exact failure tests produced by the new work.  They are intentionally short: the
point is to prevent future effort from rebuilding routes already falsified by a one-line control.

\subsection{Same-prime congruence transfer fails at \texorpdfstring{$x=2$}{x=2}}

Given an integer pair $(M,A)$ and a prime $p\nmid6MA$, consider
\[
  \phi_p(r,s)=2^rM^s\pmod p,
  \qquad
  \psi_p(r,s)=3^rA^s\pmod p.
\]
Over the reals, a common exponent identifies both expressions with powers having exponent
$r+sx$.  It is tempting to ask for the finite relation transfer
\begin{equation}
  \ker\phi_p\subseteq\ker\psi_p.
  \label{eq:same-prime-kernel-inclusion}
\end{equation}
This is not a necessary condition for the conjecture's valid solutions.

\begin{proposition}[Integral-control counterexample to same-prime transfer]
\label{prop:mod7-control}
For $x=2$, $(M,A)=(4,9)$, and $p=7$, the vector $(3,0)$ belongs to
$\ker\phi_7$ but not to $\ker\psi_7$.
\end{proposition}

\begin{proof}
One has $2^3\equiv1\pmod7$, whereas $3^3\equiv27\equiv6\pmod7$.
\end{proof}

Thus no argument may infer \eqref{eq:same-prime-kernel-inclusion} directly from the real common
exponent, whether at one prime or uniformly at every prime.  A viable local-to-global theorem must
use a different-prime modulus chosen to dominate the target order, an almost-everywhere support
hypothesis proved by independent arithmetic input, or a richer joint Kummer condition.  The real
identity alone supplies none of these.

\subsection{A compact failure matrix}

\begin{longtable}{L{0.23\textwidth}L{0.30\textwidth}L{0.36\textwidth}}
\toprule
Route & Exact obstruction & What remains viable \\
\midrule
\endfirsthead
\toprule
Route & Exact obstruction & What remains viable \\
\midrule
\endhead
Fixed finite monomial library & Theorem~\ref{thm:support-conditioning}: every fixed support becomes extinct beyond an effective block. & Supports whose size, frequency diameter, or near-resonance quality grows with $L$. \\
Pure toric scaling & Theorem~\ref{thm:toric-semi-invariant}: polynomial semi-invariants of $(X,Y)\mapsto(2X,3Y)$ are monomials. & A nondiagonal operator, a quotient of an infinite module, or a tensor construction. \\
Fixed-endpoint differential logs & Theorem~\ref{thm:monodromy-independence} and \eqref{eq:differential-integral-collapse}: specialization creates even the control relations. & A relation-faithful Kummer or motivic lift in which multiplication is functional before evaluation. \\
Naive Mahler pullback & Equation~\eqref{eq:root-alphabet}: repeated pullback creates an infinite root alphabet. & A finite invariant quotient of the root tree or a different operator family. \\
Dense near-curve support & Equation~\eqref{eq:dense-capacity}: the predicted total support capacity is never $o(L)$. & Sparse support-restricted matrices and certificate covers. \\
Fixed-arity additive energy & Equation~\eqref{eq:energy-shape-wall}: its natural growing-arity continuation costs $\exp(\Theta(n\log R))$. & Boundary-conditioned fixed-target classification of algebraic component signatures. \\
Same-prime congruence transfer & Proposition~\ref{prop:mod7-control}: it fails for $x=2$ modulo $7$. & Output-specific Kummer information or a modulus-changing continuity theorem. \\
\bottomrule
\end{longtable}

Three common mistakes are thereby separated.

\begin{enumerate}
\item \textbf{A search algorithm is not a proof criterion.}  Faster enumeration of a block is
valuable, but the proof objective is the zero-capacity inequality of
Theorem~\ref{thm:cover-capacity}.
\item \textbf{A functional relation theorem is not a value theorem.}  One must audit which
relations are created at the specialization point; integral controls are the fastest test.
\item \textbf{A local shadow is not an archimedean transfer.}  The equality of real logarithmic
ratios does not by itself impose compatible finite kernels.
\end{enumerate}

\begin{statusbox}{Proposition~\ref{prop:mod7-control} is \statusproved.  Every row of the failure
matrix points to a theorem or exact calculation in this report.  The surviving alternatives are
research directions, not claims of existence.}
\end{statusbox}

\section{Formalization architecture and certificate checking}
\label{sec:formalization}

The new arguments divide naturally into a small unconditional core, external theorem interfaces,
and computational certificates.  This is preferable to formalizing one monolithic proof attempt:
failed searches remain untrusted, while every accepted row has a short kernel-checkable meaning.

\subsection{Dependency map}

\begin{longtable}{L{0.23\textwidth}L{0.30\textwidth}L{0.36\textwidth}}
\toprule
Component & Formal content & Dependencies \\
\midrule
\endfirsthead
\toprule
Component & Formal content & Dependencies \\
\midrule
\endhead
Rational exponent lemma & If $2^{a/b}\in\Z$ in lowest terms, then $b=1$. & Integer valuations only. \\
Block pressure & Construction $n_k=L-\lfloor kx\rfloor$ and distinctness of $\{kx\}$. & Ordered fields, floor arithmetic, irrationality. \\
Dyadic normalization & $(2^Lz)^{\tht}=3^Lz^{\tht}$. & Real logarithm identities. \\
Integer locking & An integer of absolute value $<1$ is zero. & Discrete ordered ring. \\
Exponential zero budget & A $T$-term real exponential sum has at most $T-1$ zeros. & Rolle's theorem and induction. \\
Curve zero budget & Distinctness of $i+j\tht$ and substitution $z=e^t$. & Irrationality of $\tht$ plus the preceding item. \\
Cover capacity & Assign each exact point to a covering certificate and sum zero budgets. & Finite sets and cardinal inequalities. \\
Support conditioning & Explicit inverse-Vandermonde row bound. & Polynomial coefficient bounds, mean-value theorem, finite matrices. \\
Toric semi-invariants & Compare monomial coefficients under $P(2X,3Y)=cP(X,Y)$. & Unique factorization of monomials. \\
Singular matrix plane & A two-dimensional all-singular subspace with surjective first row has a fixed row multiplier. & $2\times2$ determinants and linear algebra. \\
Monodromy independence & Minimal polynomial relation plus translation monodromy. & Analytic continuation around distinct logarithmic branch points. \\
Same-prime control failure & $2^3\equiv1$, $3^3\equiv6\pmod7$. & Decidable modular arithmetic. \\
\bottomrule
\end{longtable}

The following results should enter as named interfaces rather than be reproved in the first
formalization pass:

\begin{description}
\item[Baker interface.] A polynomial lower bound for nonzero $a+b\tht$.
\item[Near-curve interface.] The Brisebarre--Hanrot rounded Chebyshev norm theorem.
\item[Carlitz interface.] Papanikolas's relation-ideal theorem for algebraic Carlitz logarithms.
\item[Additive interface.] Harrison's fixed-arity non-permutation solution bound.
\end{description}

Each interface is used in one sharply delimited theorem.  Removing it leaves the rest of the file
executable and exposes the exact unformalized assumption.

\subsection{Suggested Lean theorem signatures}

The names below are schematic but sufficiently precise to guide implementation.

\begin{lstlisting}[style=pseudostyle]
theorem rational_exponent_integral
    {a b M : Int} (hb : 0 < b) (hcop : IsCoprime a b)
    (hpow : (2 : Real) ^ ((a : Real) / b) = M) : b = 1

theorem block_pressure
    (hx : Irrational x) (hsol2 : (2 : Real)^x = M)
    (hsol3 : (3 : Real)^x = A) :
    card (blockPoints x L) >= ceil ((L + 1) / x)

theorem exp_sum_zero_count
    (hfreq : StrictMono lambda) (hcoeff : forall i, coeff i != 0) :
    card (zerosOn (fun t => sum i, coeff i * exp (lambda i * t)) I)
      <= T - 1

theorem curve_polynomial_zero_count
    (htheta : Irrational theta) (hP : P != 0) :
    card {z in I | P(z, z^theta) = 0} <= supportCard P - 1

theorem certified_cover_capacity
    (hcover : Covers unitInterval certs)
    (hsubunit : forall C in certs, IsSubunitCertificate L C) :
    card (integerCurvePoints L) <= sum C in certs, supportCard C.poly - 1

theorem support_conditioning_barrier
    (hsubunit : supNormOnBlock theta L P < 1) :
    L < conditioningBudget theta P.support

theorem toric_semiInvariant_is_monomial
    (h : renameScale 2 3 P = c * P) : IsMonomial P

theorem singular_plane_fixed_multiplier
    (hsing : forall X in W, det X = 0)
    (hsurj : Surjective firstRowOnW) :
    exists q, forall X in W, secondRow X = q * firstRow X
\end{lstlisting}

The block-pressure theorem should represent points by their real exponent $s_k$ first and only
then map to integer pairs.  This avoids proving accidental injectivity through transcendental
coordinates.  The curve-zero theorem should normalize the support into an ordered list of distinct
frequencies and reuse one generic exponential-sum theorem.

\subsection{A minimal certificate object}

An untrusted search program should emit records of the following logical form.

\begin{lstlisting}[style=pseudostyle]
structure NearCurveCertificate where
  block       : Nat
  left right  : Rat
  hInterval   : left < right
  support     : Array (Nat x Nat)
  coeff       : Array Int
  N1 N2       : Nat
  rho1 rho2   : Rat
  thetaBox    : Rat x Rat
  chebBox     : Array (Rat x Rat)
  tailBox     : Array (Rat x Rat)
  roundedA    : Matrix Int
  scaleExp    : Nat
  normSqUpper : Rat
  hNorm       : normSqUpper < 1 / (support.size + N1*N2)^2
  hAnalytic   : CertifiedAnalyticEnvelope
\end{lstlisting}

The checker performs no lattice reduction.  It verifies array dimensions, support uniqueness,
rational interval containment, directed-rounding enclosures, the exact squared-norm inequality,
and the analytic tail theorem.  The resulting proposition is
\[
  \forall z\in[a,b],\quad
  |P(2^Lz,3^Lz^{\tht})|<1.
\]
A cover certificate is a finite array of these records plus an exact rational proof that the
intervals cover $[1,2]$.  Its final output is the integer
\[
  \operatorname{Cap}=\sum_C(|\supp P_C|-1),
\]
which is compared directly with the orbit-pressure lower bound.

\subsection{Trusted and untrusted layers}

\begin{longtable}{L{0.22\textwidth}L{0.70\textwidth}}
\toprule
Layer & Responsibility \\
\midrule
\endfirsthead
\toprule
Layer & Responsibility \\
\midrule
\endhead
Untrusted search & Support enumeration, Chebyshev sampling, precision choice, LLL or BKZ, interval subdivision, and candidate ranking. \\
Exact checker & Integer matrix arithmetic, rational enclosures, norm inequalities, support capacity, and interval cover. \\
Analysis library & Bernstein ellipse bounds, Chebyshev tails, exponential zero theorem, and continuity facts. \\
External theorem interface & Baker separation and, in later branches, Papanikolas or Harrison. \\
Final kernel theorem & If a certified cover has capacity smaller than the forced counterexample point count, then no noninteger solution exists. \\
\bottomrule
\end{longtable}

This architecture permits formal progress before the analytic search succeeds.  In particular,
the finite checker, block-pressure theorem, zero-capacity theorem, and the degree-three pilot
certificate format can be completed independently.

\begin{statusbox}{The dependency architecture is a formalization plan.  The theorem statements
mirror results already proved in this report, but no Lean source is claimed to have been compiled.
External theorems are explicitly isolated.}
\end{statusbox}

\section{Prioritized research program}
\label{sec:program}

The report has reduced several broad ideas to executable tests.  The priorities below are ordered
by the ratio of mathematical payoff to missing infrastructure, not by conceptual prestige.

\subsection{Priority A: certified sparse near-curve search}

This is the most direct route because every successful record has unconditional proof content.
The first implementation should follow the loop below.

\begin{enumerate}
\item \textbf{Enumerate support families.}  Use total degree
$d\asymp\sqrt{L/\log L}$ and support size at most $CL/\log L$.  Include anisotropic lower sets,
near-resonant shells around convergents to $\tht$, and a small low-frequency core.
\item \textbf{Apply the conditioning prefilter.}  Compute $\mathfrak B(S)$ exactly enough to
reject every support with $\mathfrak B(S)\le L$.  This costs essentially nothing compared with
lattice reduction and prevents searches forbidden by Theorem~\ref{thm:support-conditioning}.
\item \textbf{Build the support-restricted Chebyshev matrix.}  Use the exact dyadic row scales
$2^{Li}3^{Lj}$, cached analytic envelopes for the interval shape, and directed rounding.
\item \textbf{Reduce and certify.}  LLL or BKZ proposes rows.  The independent checker verifies
\eqref{eq:sparse-BH-norm}.  A row that is small only on sample nodes is discarded.
\item \textbf{Subdivide only on failure.}  If no row certifies an interval, bisect it.  Charge the
actual support of every accepted row and monitor
\[
  \Gamma_L:=\frac1L\sum_C(|\supp P_C|-1).
\]
The asymptotic goal is $\Gamma_L\to0$, not merely a fast running time.
\end{enumerate}

The first milestones are deliberately finite:

\begin{description}
\item[A0.] Convert the degree-three rows for $L\le5$ into rigorous interval certificates or prove
that those particular rows fail between grid points.
\item[A1.] Reproduce the dense Brisebarre--Hanrot certificate on a small block with an independent
checker, then delete rows to verify Proposition~\ref{prop:sparse-BH} computationally.
\item[A2.] For $10\le L\le100$, record the best certified values of support count, degree,
conditioning budget, interval count, and $\Gamma_L$.
\item[A3.] Search for a family whose empirical capacity grows more slowly than $L$ and then replace
the observed pattern by a uniform analytic construction.
\end{description}

A negative outcome is still informative.  If the minimal certified capacity stabilizes above a
positive multiple of $L$ while support budgets approach the conditioning floor, the one-polynomial
cover strategy is likely at its natural limit and the search should move to operator or additive
structure rather than larger blind lattices.

\subsection{Priority B: analytic design of near-resonant supports}

The conditioning budget shows why convergents to $\tht$ are relevant.  If
$p/q$ approximates $\tht$, frequency pairs separated by $(p,-q)$ have gap
$|p-q\tht|$, increasing $\mathfrak B(S)$.  A pure chain in this direction is not sufficient: it
reduces to a gap-power construction already ruled out in Report III.  The viable design is a
\emph{two-dimensional} support with a small core and several translated near-resonant pairs.

A concrete family to analyze is
\begin{equation}
  S(d,p,q,h):=
  \{(i,j):0\le i+j\le d_0\}
  \cup
  \bigcup_{r=1}^{h}\{(i_r,j_r),(i_r+p,j_r-q)\},
  \label{eq:core-shell-support}
\end{equation}
with all coordinates nonnegative and total degree at most $d$.  The core controls low-order
Chebyshev behavior; the shell supplies small frequency gaps without collapsing the whole
polynomial to one variable.  The research tasks are:

\begin{enumerate}
\item derive a determinant or singular-value bound sharper than the worst-case
Theorem~\ref{thm:support-conditioning} for \eqref{eq:core-shell-support};
\item estimate the Brisebarre--Hanrot tail weights monomial by monomial;
\item determine whether $h\asymp L/\log L$ can be supported at
$d\asymp\sqrt{L/\log L}$; and
\item prove that any successful short vector is not forced into a line-supported factor.
\end{enumerate}

The success criterion is an explicit support family and coefficient-selection theorem, not a list
of isolated LLL rows.

\subsection{Priority C: relation-faithful operator or motivic lift}

The Carlitz theorem proves that the desired multiplier rigidity is natural in a category with a
strong difference operator.  The two ordinary-log audits identify the specifications:

\begin{enumerate}
\item multiplication of algebraic parameters must produce functional linear relations;
\item the degree-two determinant must lie in a finite tensor or quotient module;
\item the operator must not diagonalize monomials as the toric scaling does;
\item specialization must preserve the degree-two relation ideal while allowing the known linear
relations of integral controls.
\end{enumerate}

The next precise calculation is to express
\[
  \log M\otimes\log3-\log A\otimes\log2
\]
as a period tensor of the direct sum of the Kummer $1$-motives attached to
$2,3,M,A$, and compute the smallest Tannakian subquotient containing it.  One should then ask for a
\emph{single-tensor} period statement on that subquotient, rather than invoke the full
Grothendieck--Andre or Schanuel conjecture.  The characteristic-$p$ proof in
Theorem~\ref{thm:Carlitz-multiplier} supplies the model answer that the relation ideal should have.

A proposed special-function family must be tested immediately on $(M,A)=(4,9)$ and $(8,27)$.
If the value-linear relations do not lift functionally, the family is rejected before any
algebraic-independence argument is attempted.

\subsection{Priority D: growing-arity algebraic-part classification}

For Harrison's route, improving the exponent of the transcendental point count is not the first
task.  The carry points already lie in algebraic parts.  The useful program is:

\begin{enumerate}
\item parameterize character-balanced semi-algebraic components for equal-sum equations on the
triangular log-grid;
\item encode a component by a finite partition and a lattice of multiplicative characters;
\item bound the number of component signatures meeting a fixed boundary target when arity is
$O(R)$; and
\item isolate the signature operation corresponding to interchanging the two primitive carry
directions.
\end{enumerate}

A theorem with total count $|\cG_R|^{O(n)}$ should be regarded as a failure, because
\eqref{eq:energy-shape-wall} remains.  The target is exponential in $n$ with a controlled base, or
a structural statement that forbids two independent integer carry directions for irrational
$c$.

\subsection{Priority E: local arithmetic only with an explicit transfer lemma}

The existing Kummer densities and support theorems are powerful once a local hypothesis is
available.  Proposition~\ref{prop:mod7-control} shows that the real common exponent does not provide
the simplest one.  New local work should therefore begin by stating an archimedean-to-local
transfer lemma in full, including its quantifiers over primes and depth.  Unless that lemma
survives all integral controls and yields more than a same-prime kernel inclusion, no further
Chebotarev computation is warranted.

\subsection{Decision table}

\begin{longtable}{L{0.18\textwidth}L{0.34\textwidth}L{0.39\textwidth}}
\toprule
Track & Continue when & Stop or redesign when \\
\midrule
\endfirsthead
\toprule
Track & Continue when & Stop or redesign when \\
\midrule
\endhead
Sparse certificates & Certified $\Gamma_L$ decreases and surviving supports have $\mathfrak B(S)>L$ by a growing margin. & Only sample-small rows appear, or certified capacity remains $\gg L$ near the conditioning floor. \\
Near-resonant design & A two-dimensional core-shell support improves certified singular values without line factorization. & Every good vector reduces to a one-variable gap power. \\
Operator lift & Integral-control linear relations lift before specialization and the determinant lies in a finite tensor quotient. & The family is diagonal, has an uncontrolled infinite alphabet, or creates control relations only at the value. \\
Additive classification & The number of boundary component signatures has exponential rate independent of $\log R$. & The estimate retains a factor $R^{cn}$ or counts only the transcendental part. \\
Local arithmetic & A proved transfer supplies compatible Kummer data for a positive-density or growing-depth prime family. & The hypothesis is simply inferred from equality of real logarithmic ratios. \\
\bottomrule
\end{longtable}

\section{Conclusion}
\label{sec:conclusion}

No proof of the \AEconj{} conjecture has been obtained.  The new progress is nevertheless concrete and
substantially narrows three routes.

On the auxiliary-polynomial side, the block normalization is exact and scale-free:
$f(z)=z^{\tht}$, $u=2^L$, $v=3^L$.  A counterexample would place linearly many exact points in
every block, while a $T$-term polynomial can vanish at at most $T-1$ positive curve points.  This
yields the first proof criterion of the report:
\[
  \text{a certified cover of total support }o(L)\text{ proves the conjecture.}
\]
The unconditional inverse-Vandermonde argument shows that a global certificate must have degree
at least $\gg\sqrt{L/\log L}$; the counterexample pressure requires $\gg\sqrt L$.  The remaining
construction window is logarithmic.  Brisebarre--Hanrot's 2026 algorithm can be restricted
rigorously to arbitrary sparse supports, turning this criterion into an executable search and
formal-certificate program.

On the transcendence side, Papanikolas's theorem gives an exact positive-characteristic model:
two independent algebraic Carlitz logarithms cannot be multiplied by a common scalar and remain
algebraic Carlitz logarithms unless the scalar is rational over the function field.  Linear algebra
shows that the ordinary conjecture would follow from the much narrower statement that one
quadratic determinant relation among four logarithms is generated by their rational linear
relations.  Pure toric scaling cannot provide the operator.  The naive Mahler family has an
infinite root alphabet, while the naive differential family has the opposite defect: its germs are
algebraically independent, but specialization creates even $\log4=2\log2$.  A successful lift must
therefore be relation-faithful before specialization.

On the additive side, Harrison's fixed-arity theorem gives polynomial bounds on the triangular
smooth grid.  Passing from total energy to a target fiber and then allowing arity proportional to
$R$ introduces an unavoidable $\exp(\Theta(n\log R))$ wall, far above the inherited carry entropy.
The required theorem must condition on the smooth boundary target and classify the
character-balanced algebraic components; improving only the transcendental point count cannot
succeed.

The most immediate opportunity is the sparse near-curve program because its output is
self-certifying and its failure modes are quantitative.  The deepest conceptual opportunity is the
single-tensor operator target suggested by the Carlitz analogue.  The additive route becomes
competitive only after its algebraic parts are counted with growing-arity, rank-sensitive
constants.  These three directions are independent enough that progress in any one would change
the status of the conjecture, and precise enough that future work can be judged against explicit
thresholds rather than qualitative plausibility.

\begin{resultbox}
\textbf{Sharp remaining bottlenecks.}
\begin{enumerate}
\item Construct a certified dyadic cover with total monomial capacity $o(L)$.
\item Prove degree-two relation descent for the determinant tensor of
$\log2,\log3,\log M,\log A$.
\item Prove a boundary-conditioned growing-arity classification whose exponential rate beats the
mixed carry rate.
\end{enumerate}
Any one of these closes the conjecture; none is claimed here.
\end{resultbox}

\appendix

\section{Logical dependency map and certificate formats}
\label{app:dependencies}

This appendix separates the report's unconditional mathematics from the three conditional closure
routes.  Its purpose is practical: a proof assistant, an independent reader, or a computational
team should be able to replace one module without silently changing another.

\subsection{The dyadic route as a short implication chain}

For a fixed real $x$ satisfying $M=2^x\in\Z$ and $A=3^x\in\Z$, the dyadic argument has the
following dependency graph:
\[
\begin{array}{c}
 x\notin\Z
 \Longrightarrow x\notin\Q
 \Longrightarrow \#\mathcal B_L\ge \ceil{(L+1)/x}
 \\
 \Downarrow
 \\
 \text{subunit polynomial on a covering interval}
 \Longrightarrow \text{all exact integer points there are zeros}
 \\
 \Downarrow
 \\
 P(X,X^\tht)=\sum c_{ij}X^{i+j\tht}
 \Longrightarrow \#\{P=0\}\le T(P)-1
 \\
 \Downarrow
 \\
 \operatorname{Cap}_L\ge \ceil{(L+1)/x}.
\end{array}
\]
Every arrow in this diagram is proved in Sections~\ref{sec:setup}--\ref{sec:fewnomial}.  No
transcendence theorem beyond the elementary irrationality of $\log_2 3$ is used.  The sole missing
input for a proof is the construction of covers whose capacity is $o(L)$.

A machine-checkable global certificate for block $L$ may therefore consist of records
\begin{equation}
 \mathfrak C_{L,r}=(a_r,b_r,S_r,(c_{ij})_{(i,j)\in S_r},\varepsilon_r),
 \label{eq:cover-record}
\end{equation}
where all interval endpoints and error bounds are rational, together with proofs of
\begin{align}
 [1,2]&\subseteq\bigcup_r[a_r,b_r],
 \label{eq:cover-check-a}\\
 \sup_{z\in[a_r,b_r]}
 \left|\sum_{(i,j)\in S_r}c_{ij}2^{Li}3^{Lj}z^{i+j\tht}\right|
 &\le1-\varepsilon_r,
 \qquad \varepsilon_r>0,
 \label{eq:cover-check-b}\\
 \sum_r(|S_r|-1)&<L/C
 \label{eq:cover-check-c}
\end{align}
for a sequence of constants $C\to\infty$.  The analytic inequality may be certified by rational
Chebyshev enclosures, Bernstein ellipses, Taylor models, or any other validated method.  The proof
kernel need not trust the lattice reduction that proposed the coefficients.

\subsection{A support-conditioning prefilter}

Before attempting interval certification, a support $S$ can be rejected whenever
$L\ge\mathfrak B(S)$.  A rational implementation should not evaluate $\tht$ as an unproved
floating-point constant.  Instead, it may carry directed rational intervals
\[
 \underline\tht<\frac{\log3}{\log2}<\overline\tht
\]
obtained from certified logarithm bounds.  From these one obtains:

\begin{enumerate}
\item an upper bound for every frequency $i+j\tht$ and hence for $\Lambda(S)$;
\item a lower bound for every nonzero difference
$(i-i')+(j-j')\tht$ and hence for $\delta(S)$;
\item an outward-rounded upper bound for $\mathfrak B(S)$.
\end{enumerate}

For large degree, Baker--W\"ustholz or Matveev supplies a symbolic lower bound for the third item.
For moderate supports, direct certified evaluation of the finitely many frequency differences is
usually much sharper.  Thus the theoretical degree barrier and the practical pruning test use the
same theorem but should not use the same numerical constants.

\subsection{The operator route as relation-ideal descent}

The function-field part has an equally short logical shape.  For Carlitz logarithms,
Papanikolas gives
\[
 k\text{-linear independence}
 \Longrightarrow
 \overline{k}\text{-algebraic independence}.
\]
Choosing a basis converts this into the statement that every polynomial relation is generated by
linear relations.  Applying that statement to the determinant
$X_1X_4-X_2X_3$ and invoking Lemma~\ref{lem:singular-subspace} forces the common multiplier to lie
in $k$.

For ordinary logarithms, no full algebraic-independence theorem is required.  A certificate of the
needed relation descent would be an identity
\begin{equation}
 X_1X_4-X_2X_3=\sum_{r=1}^m L_r(X_1,X_2,X_3,X_4)Q_r(X_1,X_2,X_3,X_4),
 \label{eq:relation-ideal-certificate}
\end{equation}
where the $L_r$ are rational linear forms known to vanish at
$(\log2,\log3,\log M,\log A)$ and the $Q_r$ are linear forms over $\overline\Q$.  Since the left
side has degree two, only the degree-one part of each $Q_r$ matters.  Hence a successful operator
theorem ultimately produces a finite matrix identity, even if its proof uses a much larger
functional system.

The differential audit in Section~\ref{sec:differential-audit} imposes a necessary test on every
proposed lift: it must reproduce the valid integral-control relation
\[
 \log4-2\log2=0
\]
at the functional level, or prove by a regular-specialization theorem that this relation descends
from the chosen functional relation ideal.  Functional algebraic independence by itself is not an
advantage if specialization creates new relations.

\subsection{The additive route and the required theorem shape}

Let $\cG_R=\{2^a3^b:a+b\le R\}$.  A theorem useful for the conjecture must distinguish three
quantifiers that are harmless at fixed arity but decisive when $n\asymp R$:

\begin{enumerate}
\item the exponent $x$ is a fixed irrational number satisfying the two-base integrality
hypothesis;
\item the arity $n$ grows linearly with $R$;
\item the target is not arbitrary but belongs to the smooth boundary grid
$\cG_{O(R)}^{[x]}$ and is accompanied by character-balance data inherited from the two generators.
\end{enumerate}

A prototype conclusion strong enough to interact with the carry lower bound would have the form
\begin{equation}
 \max_{T\in\cG_{cR}^{[x]}}
 \#\{(g_1,\ldots,g_n)\in\cG_R^n:g_1^x+\cdots+g_n^x=T\}
 \le \exp\bigl((\sigma+o(1))n\bigr),
 \label{eq:additive-prototype}
\end{equation}
with $\sigma$ strictly smaller than the independently certified carry entropy.  Any estimate with
an uncontrolled factor $R^{cn}$ has logarithm $cn\log R$ and cannot meet this threshold.
Estimate~\eqref{eq:additive-prototype} is not asserted; it records the exact asymptotic shape that a
future higher-rank theorem must achieve.

\subsection{Closure conditions and independence of the routes}

The three closure modules are logically independent:

\begin{longtable}{L{0.25\textwidth}L{0.66\textwidth}}
\toprule
Route & Sufficient new input \\
\midrule
\endfirsthead
\toprule
Route & Sufficient new input \\
\midrule
\endhead
Dyadic auxiliary polynomials & Certified covers with
$\sum_r(|S_{L,r}|-1)=o(L)$ along an unbounded sequence of block indices. \\
Operator/transcendence & Degree-two relation descent for the determinant tensor of the four
ordinary logarithms, uniformly for every putative solution. \\
Additive classification & A boundary-conditioned growing-arity fiber upper bound whose
exponential rate is strictly below the imported carry-fiber lower rate. \\
\bottomrule
\end{longtable}

A proof from the first route would not require the operator or additive targets.  A proof from the
second would not require any computational certificate.  A proof from the third would use the
carry construction from Report III but not the near-curve algorithm.  This modularity is useful:
a failed experiment in one track does not weaken the exact reductions in the others.

\section{Reproducibility script}
\label{app:script}

The following Python script reproduces the numerical constants in
Tables~\ref{tab:cheb}, \ref{tab:lll-pilot}, and \ref{tab:BH-exponents}, including the explicit
polynomial~\eqref{eq:pilot-polynomial}.  It uses arbitrary-precision evaluation for reconnaissance
and exact integer arithmetic for LLL.  As emphasized in the main text, the sampled supremum is not
an interval certificate.

\begin{lstlisting}[style=pythonstyle,caption={Reproduction of the numerical tables and the degree-three LLL pilot.}]
"""Reproduce the numerical data in Continuation Report IV.

Requirements:
    Python 3.11 or later
    mpmath
    sympy

The LLL lattice is exact.  The final supremum is sampled on 2001 points
and is therefore diagnostic, not a proof of a uniform inequality.
"""

from __future__ import annotations

from dataclasses import dataclass
from typing import Iterable

import mpmath as mp
from sympy import Matrix

mp.mp.dps = 100
THETA = mp.log(3) / mp.log(2)


def chebyshev_table() -> None:
    """Print the rho=5 error and conservative block-scale table."""
    constant = mp.power(mp.mpf("2.8"), THETA)
    print("Chebyshev diagnostics (rho=5)")
    print("degree       E_n                    B")
    for degree in (8, 12, 20, 28, 36, 48):
        error = constant * mp.power(5, -degree)
        block_scale = mp.power(1 / (8 * error), 1 / THETA)
        print(
            f"{degree:>6d}  "
            f"{mp.nstr(error, 12):>20s}  "
            f"{mp.nstr(block_scale, 12):>20s}"
        )


def bh_exponent_table() -> None:
    """Print the dense-support exponents in the dyadic variables."""
    print("\nBrisebarre--Hanrot exponent translation")
    print("d       N_d            a_d          kappa_d")
    for degree in (4, 6, 8, 10, 12, 16, 20, 24, 32, 40, 64):
        dimension = (degree + 1) * (degree + 2) // 2
        a_d = mp.mpf(2 * degree) / (3 * dimension)
        kappa_d = (1 + THETA) * a_d
        print(
            f"{degree:>2d}  {dimension:>8d}  "
            f"{float(a_d):>13.6f}  {float(kappa_d):>13.6f}"
        )


@dataclass(frozen=True)
class PilotResult:
    block: int
    sampled_supremum: mp.mpf
    coefficients: tuple[int, ...]


def total_degree_monomials(degree: int) -> list[tuple[int, int]]:
    """Order monomials by Y-degree, then X-degree."""
    return [
        (i, j)
        for j in range(degree + 1)
        for i in range(degree - j + 1)
    ]


def evaluate_polynomial(
    coefficients: Iterable[int],
    monomials: list[tuple[int, int]],
    block: int,
    z: mp.mpf,
) -> mp.mpf:
    """Evaluate P(2^L z, 3^L z^theta) at high precision."""
    x_value = mp.power(2, block) * z
    y_value = mp.power(3, block) * mp.power(z, THETA)
    total = mp.mpf("0")
    for coefficient, (i, j) in zip(coefficients, monomials):
        total += (
            coefficient
            * mp.power(x_value, i)
            * mp.power(y_value, j)
        )
    return total


def lll_pilot(block: int, degree: int = 3) -> PilotResult:
    """Run the exact sample lattice used in the report."""
    monomials = total_degree_monomials(degree)
    dimension = len(monomials)
    coefficient_scale = 10**8
    evaluation_scale = 10**30

    nodes = [
        mp.mpf("1.5")
        + mp.mpf("0.5")
        * mp.cos((k + mp.mpf("0.5")) * mp.pi / dimension)
        for k in range(dimension)
    ]

    rows: list[list[int]] = []
    for row_index, (i, j) in enumerate(monomials):
        row = [0] * dimension
        row[row_index] = coefficient_scale
        for z in nodes:
            value = (
                mp.power(2, block * i)
                * mp.power(3, block * j)
                * mp.power(z, i + j * THETA)
            )
            row.append(int(mp.nint(evaluation_scale * value)))
        rows.append(row)

    reduced = Matrix(rows).lll()
    best_supremum: mp.mpf | None = None
    best_coefficients: tuple[int, ...] | None = None

    for row_index in range(reduced.rows):
        first_coordinates = [
            int(reduced[row_index, column])
            for column in range(dimension)
        ]
        if any(value % coefficient_scale for value in first_coordinates):
            raise ArithmeticError("LLL row left the coefficient sublattice")
        coefficients = tuple(
            value // coefficient_scale for value in first_coordinates
        )

        # Exclude the trivial constant polynomial, whose sampled norm is one.
        if all(value == 0 for value in coefficients[1:]):
            continue

        sampled_supremum = mp.mpf("0")
        for grid_index in range(2001):
            z = mp.mpf(1) + mp.mpf(grid_index) / 2000
            sampled_supremum = max(
                sampled_supremum,
                abs(
                    evaluate_polynomial(
                        coefficients,
                        monomials,
                        block,
                        z,
                    )
                ),
            )

        if (
            best_supremum is None
            or sampled_supremum < best_supremum
        ):
            best_supremum = sampled_supremum
            best_coefficients = coefficients

    if best_supremum is None or best_coefficients is None:
        raise RuntimeError("No nonconstant reduced row was found")

    return PilotResult(block, best_supremum, best_coefficients)


def format_polynomial(
    coefficients: tuple[int, ...],
    monomials: list[tuple[int, int]],
) -> str:
    """Return a compact machine-readable polynomial string."""
    terms = []
    for coefficient, (i, j) in zip(coefficients, monomials):
        if coefficient == 0:
            continue
        terms.append(f"({coefficient})*X^{i}*Y^{j}")
    return " + ".join(terms)


def pilot_table() -> None:
    print("\nDegree-three LLL pilot")
    monomials = total_degree_monomials(3)
    for block in range(1, 7):
        result = lll_pilot(block)
        print(
            f"L={block}: "
            f"sampled sup={mp.nstr(result.sampled_supremum, 15)}"
        )
        if block == 5:
            print("L=5 coefficients:", result.coefficients)
            print("L=5 polynomial:")
            print(format_polynomial(result.coefficients, monomials))


if __name__ == "__main__":
    print("theta =", mp.nstr(THETA, 50))
    chebyshev_table()
    bh_exponent_table()
    pilot_table()
\end{lstlisting}

On the versions used for this report, the six sampled suprema are
\[
  6.9932152582\times10^{-5},\quad
  7.9924631253\times10^{-4},\quad
  6.3015062341\times10^{-3},
\]
\[
  5.8264974254\times10^{-2},\quad
  5.5874173100\times10^{-1},\quad
  4.0991943672.
\]
Exact agreement of the coefficient tuple for $L=5$ is the most useful implementation check,
because it tests the monomial ordering, lattice orientation, and scaling simultaneously.

\section{Literature and search audit}
\label{app:literature}

The literature search was intentionally narrow: only results that alter one of the report's three
quantitative routes were imported.  The following table records the role of each recent source and
what is \emph{not} inferred from it.

\begin{longtable}{L{0.21\textwidth}L{0.34\textwidth}L{0.36\textwidth}}
\toprule
Source & Imported content & Boundary of use \\
\midrule
\endfirsthead
\toprule
Source & Imported content & Boundary of use \\
\midrule
\endhead
Brisebarre--Hanrot (June 2026) & A rigorously checked short-vector condition yielding an integer
polynomial uniformly smaller than one near an analytic curve; asymptotic interval and degree
parameters for the dense triangular support. & The report does not assume their heuristic
coprimality/elimination step.  The sparse-support extension used here follows by deleting rows from
the same matrix and rechecking the norm implication. \\
Estienne (March 2026) & A Mahler-method proof of algebraic independence for linearly independent
Carlitz logarithms, exhibiting a finite semilinear functional system and a specialization theorem.
& No characteristic-zero analogue is claimed.  The paper is used to identify the operator and
specialization features that a new argument would need. \\
Harrison (December 2025; revised July 2026) & Fixed-arity bounds for additive relations among
irrational powers of sets lying in generalized arithmetic progressions, including logarithmic
losses depending on the ambient scale. & The constants are not uniform in arity, and the theorem
does not classify the boundary-conditioned algebraic components required here. \\
Papanikolas (2008) & Algebraic independence of $k$-linearly independent Carlitz logarithms of
algebraic points. & The multiplier-rigidity theorem is a corollary proved in this report; it is a
positive-characteristic benchmark, not a proof of the ordinary conjecture. \\
Baker--W\"ustholz; Matveev & Effective lower bounds for nonzero linear forms in logarithms of fixed
algebraic numbers. & Only the existence of an effective polynomial separation
$|a+b\tht|\gg H^{-C}$ is used; no optimized numerical constant is asserted. \\
\bottomrule
\end{longtable}

The search also checked the competitive claims in the prompt.  Those sources establish useful
context, but none contains a disclosed approach to the \AEconj{} conjecture.  In particular, no
inference about the rival team's private Lean development enters any proof or target in this
report.

\begin{thebibliography}{99}

\bibitem{AnthropicFable2026}
Anthropic,
\emph{Claude Fable 5 and Claude Mythos 5},
announcement dated 9 June 2026, with service updates dated 12 June and 1 July 2026,
\url{https://www.anthropic.com/news/claude-fable-5-mythos-5}.

\bibitem{BakerWustholz}
A.~Baker and G.~W\"ustholz,
\emph{Logarithmic forms and group varieties},
J. Reine Angew. Math. \textbf{442} (1993), 19--62,
\url{https://doi.org/10.1515/crll.1993.442.19}.

\bibitem{BrisebarreHanrot2026}
N.~Brisebarre and G.~Hanrot,
\emph{Integer points close to a transcendental curve: an algorithmic approach},
arXiv:2606.04858v1 (2026),
\url{https://arxiv.org/abs/2606.04858}.

\bibitem{DujellaRank2026}
A.~Dujella,
\emph{History of elliptic curves rank records},
updated August 2026,
\url{https://web.math.pmf.unizg.hr/~duje/tors/rankhist.html}.

\bibitem{EpochHadamard2026}
Epoch AI,
\emph{FrontierMath: Open Problems},
changelog entry of 12 August 2026 and the problem page for a Hadamard matrix of order $668$,
\url{https://epoch.ai/frontiermath/open-problems}.

\bibitem{Estienne2026}
G.~Estienne,
\emph{Mahler's method and Carlitz logarithm},
arXiv:2603.18832v2 (2026),
\url{https://arxiv.org/abs/2603.18832}.

\bibitem{Harrison2026}
J.~Harrison,
\emph{Additive relations in irrational powers},
arXiv:2512.04081v3 (2026),
\url{https://arxiv.org/abs/2512.04081}.

\bibitem{JacobianReport2026}
The Conversation,
\emph{Claude Fable 5 AI finds a tiny formula that topples an 87-year-old math conjecture},
republished by ScienceDaily, 5 August 2026,
\url{https://www.sciencedaily.com/releases/2026/08/260804034634.htm}.

\bibitem{MathWorldRank2026}
E.~W.~Weisstein,
\emph{Elliptic Curve Rank},
MathWorld--A Wolfram Resource, updated August 2026,
\url{https://mathworld.wolfram.com/EllipticCurveRank.html}.

\bibitem{Matveev2000}
E.~M.~Matveev,
\emph{An explicit lower bound for a homogeneous rational linear form in the logarithms of
algebraic numbers. II},
Izv. Math. \textbf{64} (2000), no.~6, 1217--1269,
\url{https://doi.org/10.1070/IM2000v064n06ABEH000314}.

\bibitem{Papanikolas2008}
M.~A.~Papanikolas,
\emph{Tannakian duality for Anderson--Drinfeld motives and algebraic independence of Carlitz
logarithms},
Invent. Math. \textbf{171} (2008), 123--174,
\url{https://doi.org/10.1007/s00222-007-0073-y}.

\end{thebibliography}

\end{document}
